% !TEX root = ../main.tex
\chapter{Probability Theory}\label{ch:probability}
\begin{paracol}{2}

\begin{sources}
\begin{tabularx}{\linewidth}{@{}l l L@{}}
\toprule
\textbf{Lecture} & \textbf{Speaker} & \textbf{Content}\\
\midrule
2024 L4 (from 27:40) & P.~Kempthorne & Discrete, continuous and mixed variables; CDF and probability integral transform; moments, skewness, kurtosis; normal and lognormal laws, change of variables; expected call payoff; moment-generating functions; independence, covariance, random vectors and portfolio variance; start of PCA.\\
2024 L5 (to 1:12:25) & P.~Kempthorne & PCA of a random vector; $\chi^2$, $t$ and $F$ laws; weak law of large numbers, central limit theorem; \emph{Probability Theory for Asset Pricing}: Stein's lemma, expected-utility pricing, CARA utility, CAPM. (The rest of L5 is in \cref{ch:sp1}.)\\
2013 L3 & C.~Lee & pmf/pdf, mutual vs pairwise independence; why prices are lognormal; change of variables; exponential families; MGFs and the moment problem; weak LLN (casino and poker); CLT with proof.\\
\bottomrule
\end{tabularx}

\smallskip
\textbf{Files.} Slides \emph{Probability Theory} \file{mit18_642_f24_lec04_1.pdf}; slides
\emph{Probability Theory for Asset Pricing} \file{mit18_642_f24_lec04_2.pdf}; 2013 notes
\file{MIT18_S096F13_lecnote3.pdf}. Suggested readings on the 2024 Week~3 page: J.~Lintner (1965),
\emph{The valuation of risk assets and the selection of risky investments in stock portfolios and
capital budgets}, and \emph{Security prices, risk, and maximal gains from diversification}.
\textbf{Problem sets.} 2024 PS2 Problem~2; 2024 PS3 Problems~1--2; 2013 PS2 A-1, A-2, B-1 to B-4.
The other problems of these sets (random walks, Markov chains, martingales, PCA of yields) are in
\cref{ch:sp1,ch:pca}.
\textbf{Not recorded.} Nothing: all three lectures are on video.
\end{sources}

\section*{Motivation}
Every model in this course is a probability model. This chapter collects the probability that the
rest of the notes take for granted, organised around four questions that a quant asks about a
position:
\begin{itemize}
  \item \emph{What is the distribution of a price or a return?} Normal for log-returns, lognormal
        for prices, and in reality heavier tails than either (\cref{ch03:sec-normal,ch03:sec-heavy}).
        The distribution fixes the expected payoff of an option (\cref{ch03:sec-call}).
  \item \emph{How do assets move together?} Covariance matrices, portfolio variance $\bm a\T\Sigma\bm a$,
        diversification, the multivariate normal law and its principal components
        (\cref{ch03:sec-multi,ch03:sec-mvn,ch03:sec-pca}).
  \item \emph{What survives aggregation?} Laws of large numbers and the central limit theorem explain
        why a small edge repeated many times makes money and why log-returns over long horizons look
        Gaussian (\cref{ch03:sec-limits}).
  \item \emph{What should an asset cost?} With nothing more than expectations, covariances and one
        Gaussian identity (Stein's lemma), an expected-utility argument yields the capital asset
        pricing model (\cref{ch03:sec-capm}).
\end{itemize}
A physicist can skim \cref{ch03:sec-rv,ch03:sec-moments} and the definitions of
\cref{ch03:sec-chisq}; the less standard material is the option-payoff formula, the pitfalls of
moment-generating functions, conditional expectation, Stein's lemma and the CAPM derivation.

% ------------------------------------------------------------------
\section{Random variables and distribution functions}\label{ch03:sec-rv}

Financial random variables come in three kinds \ts{24}{4}{28:47}:
\begin{itemize}
  \item \emph{discrete}: default (1) or no default (0) of a counterparty, the FOMC decision on the Fed
        funds rate, the side (buy/sell) and the share size of the next market order in AAPL;\begin{answer}{what is FOMC?}[FOMC]
The \emph{Federal Open Market Committee}: the body of the US Federal Reserve that sets the target for the
Fed funds rate (the overnight rate between banks). It meets eight times a year; each decision (cut, hold,
hike, usually by 0.25\%) is a discrete random variable for markets. ``Open market'' because the Fed
implements the decision by buying and selling bonds in the open market.
\end{answer}
  \item \emph{continuous}: the value of an asset, the waiting time until the next market order
        \ts{24}{4}{29:49};
  \item \emph{mixed}: the value of a stock in six months, which is continuous except for an atom at
        $0$ (bankruptcy) or at the buyout price (the company is taken private) \ts{24}{4}{30:55}.
\end{itemize}

\begin{definition}[Distribution of a random variable]
Let $\mathcal X$ be the set of possible values (the \emph{sample space}) of $X$.
\begin{itemize}
  \item If $X$ is discrete, its \emph{probability mass function} (pmf) is $f_X(x)=\Prob(X=x)\ge0$ with
        $\sum_{x\in\mathcal X}f_X(x)=1$.
  \item If $X$ is continuous, its \emph{probability density function} (pdf) $f_X\ge0$ satisfies
        $\Prob(X\in[x,x+\dd x])=f_X(x)\dd x$ and $\int_{\mathcal X}f_X(x)\dd x=1$.
  \item The \emph{cumulative distribution function} (CDF) is $F_X(x)=\Prob(X\le x)$
        \ts{24}{4}{31:56}. It is non-decreasing, right-continuous, with $F_X(-\infty)=0$ and
        $F_X(+\infty)=1$, and it determines the distribution.
\end{itemize}
Then $\Prob(X\in A)=\sum_{x\in A}f_X(x)$ or $\int_Af_X(x)\dd x$, and $F_X(x)=\int_{-\infty}^xf_X(u)\dd u$ in
the continuous case, so $f_X=F_X'$.
\end{definition}

\begin{remark}
Both courses identify the sample space with the set of values of $X$ (2013 notes, section~1). The
measure-theoretic set-up (a probability space $(\Omega,\mathcal F,\Prob)$ and $X$ a measurable map
$\Omega\to\R$) is not needed here; it becomes useful for martingales and filtrations in \cref{ch:sp1}.
\end{remark}

\paragraph{Independence.} $X$ and $Y$ are \emph{independent} if
$\Prob(\{X\in A\}\cap\{Y\in B\})=\Prob(X\in A)\,\Prob(Y\in B)$ for all sets $A,B$; equivalently, the joint
pmf/pdf factorises, $f_{X,Y}(x,y)=f_X(x)f_Y(y)$ \ts{13}{3}{6:16}. For $n$ variables one must distinguish
\ts{13}{3}{7:17}: $X_1,\dots,X_n$ are \emph{mutually independent} if
$\Prob(\bigcap_i\{X_i\in A_i\})=\prod_i\Prob(X_i\in A_i)$ for all $A_1,\dots,A_n$, and \emph{pairwise
independent} if every pair is independent. Mutual implies pairwise, not conversely
(\cref{ch03:ex-pairwise}). \note+{interesting new subtelty}``Independent'' always means mutually independent.

\begin{proposition}[Probability integral transform]\label{ch03:prop-pit}
Let $X$ have a continuous CDF $F$. Then $U=F(X)$ is uniform on $(0,1)$ \ts{24}{4}{32:57}. Conversely, if
$U\sim\mathrm{Uniform}(0,1)$ and $F^{-1}(u)=\inf\{x:F(x)\ge u\}$ is the quantile function of an
arbitrary CDF $F$, then $F^{-1}(U)$ has CDF $F$.
\end{proposition}
\begin{proof}
If $F$ is strictly increasing, $\Prob(F(X)\le u)=\Prob(X\le F^{-1}(u))=F(F^{-1}(u))=u$ for $u\in(0,1)$. In
general, with $q=F^{-1}(u)$ continuity gives $F(q)=u$, and $\{F(X)\ge u\}=\{X\ge q\}$ (if $X\ge q$ then
$F(X)\ge F(q)=u$; if $X<q$ then $F(X)<u$ by definition of the infimum), so
$\Prob(F(X)\ge u)=1-F(q)=1-u$. For the converse, $F^{-1}(u)\le x\iff u\le F(x)$ by right-continuity, so
$\Prob(F^{-1}(U)\le x)=\Prob(U\le F(x))=F(x)$.
\end{proof}

\begin{finance}[Uses of the probability integral transform]
\begin{itemize}
  \item \emph{Model checking} \ts{24}{4}{33:57}: if a risk model predicts that tomorrow's return has CDF
        $F_t$, the numbers $u_t=F_t(x_t)$ computed from realised returns should look like an i.i.d.\
        uniform sample. A histogram of the $u_t$ piling up near $0$ and $1$ reveals tails that are too
        thin: this is how value-at-risk and density forecasts are back-tested (\cref{ch:var}).\note+{cool way to spot the heavy tails}
  \item \emph{Simulation}: $F^{-1}(U)$ turns uniform random numbers into draws from any distribution
        (inverse-transform sampling).
  \item \emph{Copulas} (not discussed in class): applying $F_i$ to each coordinate of a random vector
        separates the marginal laws from the dependence structure.\begin{answer}{non credo di aver capito quest'ultima}[Copulas]
Take two returns $(X,Y)$ and set $U=F_X(X)$, $V=F_Y(Y)$. By \cref{ch03:prop-pit} each of $U,V$ is uniform:
all the information about the \emph{marginal} laws (fat tails, skew) has been removed. What is left, the joint
law of $(U,V)$ on the unit square, is the pure \emph{dependence}: the \emph{copula} (Latin, ``link''). Sklar's
theorem: $F_{X,Y}(x,y)=C(F_X(x),F_Y(y))$. So one can model each asset's tails and their co-movement
separately, e.g.\ fat-tailed marginals glued with a Gaussian copula (the model blamed for CDO losses in 2008).
\end{answer}
\end{itemize}
\end{finance}

% ------------------------------------------------------------------
\section{Expectation, moments and the shape of a distribution}\label{ch03:sec-moments}

\begin{definition}[Moments, skewness, kurtosis]\label{ch03:def-moments}
For a random variable $X$ (sums replace integrals in the discrete case):
\begin{itemize}
  \item mean $\mu=\E[X]=\int x f_X(x)\dd x$; $k$-th moment $m_k=\E[X^k]$;
  \item variance $\Var(X)=\E[(X-\mu)^2]=m_2-m_1^2$ and standard deviation $\sigma=\sqrt{\Var(X)}$, which has the
        units of $X$ \ts{24}{4}{35:03};
  \item with the standardised variable $Z=(X-\mu)/\sigma$ \ts{24}{4}{36:07}, the \emph{skewness}
        $\gamma=\E[Z^3]$ and the \emph{kurtosis} $\kappa=\E[Z^4]$, both dimensionless \ts{24}{4}{37:12}.
\end{itemize}
$\gamma>0$ indicates a long right tail, $\gamma<0$ a long left tail. The normal law has $\gamma=0$ and
$\kappa=3$; $X$ is \emph{leptokurtic} (fat-tailed relative to the Gaussian) if $\kappa>3$
\ts{24}{4}{38:18}. The \emph{excess kurtosis} is $\kappa-3$.
\end{definition}

\begin{proposition}[Affine transformations]\label{ch03:prop-affine}
If $Y=a+bX$ with $b\neq0$, then $\E[Y]=a+b\mu$, $\Var(Y)=b^2\sigma^2$, $\gamma_Y=\operatorname{sign}(b)\,\gamma_X$
and $\kappa_Y=\kappa_X$.
\end{proposition}
\begin{proof}
$(Y-\E[Y])/\sigma_Y=b(X-\mu)/(|b|\sigma)=\operatorname{sign}(b)\,Z_X$; take third and fourth powers.
\end{proof}

\begin{coursenote}[Erratum (in class) and a caveat on the slides]
At \ts{24}{4}{1:02:42} it is said that linear transformations do not change the skewness but might change
the kurtosis. It is the other way round (\cref{ch03:prop-affine}): kurtosis is invariant, skewness
changes sign when $b<0$ (a short position has the opposite skew of the long one). Also, slide~6 of
\file{lec04_1} lists ``$\gamma=0$: $X$ is symmetric about $\mu$''. Symmetry implies $\gamma=0$, but not
conversely: $X\in\{-3,-1,2\}$ with probabilities $0.1,\,0.5,\,0.4$ has $\E[X]=\E[X^3]=0$ and is not symmetric.
Software differs on whether ``kurtosis'' means $\kappa$ or $\kappa-3$: check before comparing numbers.
\end{coursenote}

\begin{finance}[What the moments of returns mean]
For a return, the mean is the drift and the standard deviation the \emph{volatility} (quoted
annualised). Equity-index returns are typically negatively skewed (crashes are sudden, rallies are
gradual) and daily returns of almost all assets are strongly leptokurtic. Both facts matter for risk:
a model matched only to mean and variance underestimates the probability of large losses
(\cref{ch03:sec-heavy}), and they are visible in option markets as the volatility smile (\cref{ch:volatility}).
\end{finance}\begin{answer}{can you include some graphs of: a typical probability distribution of an equity index, of a daily return of an asset. what is the valatility smile?}[Graphs]
Done in \cref{ch03:fig-returns}, with the course's own data (daily closes of the S\&P~500 stocks,
2019--2024). The equal-weighted portfolio of all 487 stocks stands in for the index.
\answerpart{Volatility smile}
Take the Black--Scholes formula, which has one unknown, $\sigma$. From the market price of each option
solve for the $\sigma$ that reproduces it: the \emph{implied volatility}. If returns were normal (BS
true), it would be the same number for every strike. It is not: plotted against the strike it is a
curve, U-shaped (a ``smile'') or, for equity indices, falling from left to right (a ``skew'', or
``smirk''), sketched in panel (c). Options far from the money are dearer than BS says, because the
market knows the tails are fat (and, for indices, that crashes go down). Full story in \cref{ch:volatility}.
\end{answer}

\begin{figure}[!htb]
\centering
\pgfplotsset{retplot/.style={width=5.1cm,height=4.8cm,tick label style={font=\scriptsize},label style={font=\small},title style={font=\small},legend style={font=\scriptsize},xlabel={daily log-return},
  scaled x ticks=false,xticklabel style={/pgf/number format/fixed},legend style={font=\tiny,fill=white,fill opacity=0.8,text opacity=1}}}
\begin{tikzpicture}
\begin{axis}[retplot,title={(a) linear scale},xmin=-0.08,xmax=0.08,ymin=0,ymax=75,xtick={-0.05,0,0.05},legend pos=north west]
  \addplot[ybar interval,fill=defcol!35,draw=defcol!60] table[x=x,y=ew] {data/fm_dailyret.dat};
  \addplot[thmcol,thick,domain=-0.08:0.08,samples=120] {exp(-(x-0.00068)^2/(2*0.01386^2))/(0.01386*sqrt(2*pi))};
  \legend{S\&P~500 (EW),normal fit}
\end{axis}
\end{tikzpicture}\hfill
\begin{tikzpicture}
\begin{axis}[retplot,title={(b) log scale: the tails},xmin=-0.14,xmax=0.12,ymode=log,ymin=0.05,ymax=3000,xtick={-0.1,0,0.1},
  legend pos=north west]
  \addplot[only marks,mark size=0.9pt,defcol] table[x=x,y=ew] {data/fm_dailyret.dat};
  \addplot[only marks,mark size=0.9pt,intcol] table[x=x,y=aapl] {data/fm_dailyret.dat};
  \addplot[thmcol,thick,domain=-0.07:0.07,samples=120] {exp(-(x-0.00068)^2/(2*0.01386^2))/(0.01386*sqrt(2*pi))};
  \legend{S\&P~500 (EW),AAPL,normal (EW)}
\end{axis}
\end{tikzpicture}\hfill
\begin{tikzpicture}
\begin{axis}[width=5.1cm,height=4.8cm,tick label style={font=\scriptsize},label style={font=\small},title style={font=\small},legend style={font=\scriptsize},title={(c) implied volatility (schematic)},
  xlabel={strike $K/S_0$},ylabel={implied vol},xmin=0.7,xmax=1.3,ymin=0.1,ymax=0.4,ytick=\empty,
  legend pos=north east]
  \addplot[thmcol,thick,domain=0.7:1.3,samples=60] {0.18+1.2*(x-1)^2};
  \addplot[defcol,thick,domain=0.7:1.3,samples=60] {0.2-0.35*(x-1)+0.8*(x-1)^2};
  \addplot[black!50,dashed,domain=0.7:1.3] {0.15};
  \legend{smile,index skew,BS (flat)}
\end{axis}
\end{tikzpicture}
\caption{Daily log-returns, 2019-01-02 to 2024-08-30, from \file{SP500_data_subset.RData} (the data of the
case study in \cref{sec:ewport}). (a) The equal-weighted S\&P~500 portfolio against the normal with the same mean
and standard deviation (1.39\% a day): a sharper peak and fatter tails; skewness $-0.91$, kurtosis $17$. (b) On a log scale
the normal is a parabola and the data are not: both the index and Apple (kurtosis $8.5$) have 11 days beyond
$4\sigma$, where the normal would expect $0.09$. The extreme left points are March 2020. (c) A sketch, not data,
of the volatility smile and of the skew typical of index options (\cref{ch:volatility}).}\label{ch03:fig-returns}
\end{figure}

% ------------------------------------------------------------------
\section{Normal and lognormal distributions}\label{ch03:sec-normal}

\subsection{The normal distribution}
\begin{definition}[Normal law]
$X\sim N(\mu,\sigma^2)$, $\mu\in\R$, $\sigma>0$, if
\[
  f(x)=\frac{1}{\sigma\sqrt{2\pi}}\exp\Bigl(-\frac{(x-\mu)^2}{2\sigma^2}\Bigr),\qquad x\in\R .
\]
$\phi$ and $\Phi$ denote the density and CDF of the standard normal $N(0,1)$; $\Phi(-x)=1-\Phi(x)$.\begin{answer}{perche' non restiamo con una notazione unica? $\phi$ = $f_normal$ e $\Phi$ = $F_normal$}[Why $\phi,\Phi$]
Because $\phi$ and $\Phi$ are universal: every paper and the Black--Scholes formula use them, so you must
recognise them. They are also a special case, not a different notation: for $Z\sim N(0,1)$, $f_Z=\phi$
and $F_Z=\Phi$; for $X\sim N(\mu,\sigma^2)$, $F_X(x)=\Phi((x-\mu)/\sigma)$.
\end{answer}
\end{definition}
The density is symmetric about $\mu$, maximal at $\mu$, with inflection points at $\mu\pm\sigma$
\ts{24}{4}{39:29}; the probabilities of lying within $1,2,3$ standard deviations are
$0.6827$, $0.9545$, $0.9973$ \ts{24}{4}{40:31}. That $f$ integrates to $1$ is the Gaussian integral
(2013 PS2 B-2, \cref{ch03:ex-b2}); see \cref{ch03:fig-normal}.\note*{insert a graph of the distribution. in general, whenever it can be useful plot graphs and functions.}

\begin{figure}[!htb]
\centering
\begin{tikzpicture}
\begin{axis}[width=0.8\linewidth,height=5cm,axis lines=left,xmin=-3.6,xmax=3.6,ymin=0,ymax=0.45,
  xtick={-3,-2,-1,0,1,2,3},xticklabels={$\mu-3\sigma$,$\mu-2\sigma$,$\mu-\sigma$,$\mu$,$\mu+\sigma$,$\mu+2\sigma$,$\mu+3\sigma$},
  ytick=\empty,tick label style={font=\scriptsize},label style={font=\small},title style={font=\small},legend style={font=\scriptsize}]
  \addplot[fill=defcol!12,draw=none,domain=-3:3,samples=80] {exp(-x^2/2)/sqrt(2*pi)} \closedcycle;
  \addplot[fill=defcol!25,draw=none,domain=-2:2,samples=80] {exp(-x^2/2)/sqrt(2*pi)} \closedcycle;
  \addplot[fill=defcol!45,draw=none,domain=-1:1,samples=80] {exp(-x^2/2)/sqrt(2*pi)} \closedcycle;
  \addplot[very thick,defcol,domain=-3.6:3.6,samples=120] {exp(-x^2/2)/sqrt(2*pi)};
  \node[font=\scriptsize] at (axis cs:0,0.12) {68.3\%};
  \node[font=\scriptsize] at (axis cs:1.5,0.03) {95.4\%};
  \node[font=\scriptsize] at (axis cs:2.55,0.025) {99.7\%};
  \fill[thmcol] (axis cs:-1,0.242) circle (1.5pt) (axis cs:1,0.242) circle (1.5pt);
  \node[font=\scriptsize,thmcol,anchor=west] at (axis cs:1.05,0.26) {inflection points};
\end{axis}
\end{tikzpicture}
\caption{The normal density $N(\mu,\sigma^2)$ with the probability of lying within $1$, $2$, $3$ standard deviations
of the mean. The inflection points are at $\mu\pm\sigma$.}\label{ch03:fig-normal}
\end{figure}

\subsection{Why prices are not normal}
Suppose daily price changes $P_n-P_{n-1}$ were i.i.d.\ normal \ts{13}{3}{10:28}. Two things go wrong
\ts{13}{3}{11:30}:
\begin{itemize}
  \item the size of a move would not depend on the price level, whereas a \$50 stock typically moves
        about five times as many dollars as a \$10 stock: it is the \emph{percentage} change that is
        roughly stationary \ts{13}{3}{12:32};
  \item $P_n$ would be normal with variance growing like $n$, so it eventually becomes negative with
        substantial probability \ts{13}{3}{19:05}. Locally (short horizons, small moves) the additive model
        is harmless.
\end{itemize}
The remedy is to model \emph{log-returns} $r_n=\ln(P_n/P_{n-1})$ as i.i.d.\ normal. Then
$\ln P_n=\ln P_0+\sum_{i\le n}r_i$ is normal (sums of independent normals are normal,
\cref{ch03:sec-mgf}), and $P_n$ is \emph{lognormal}; positivity is automatic \ts{24}{4}{41:34}. In continuous
time the same construction with Brownian increments gives geometric Brownian motion
\ts{24}{4}{42:35} (\cref{ch:stochcalc,ch:sde}).

\subsection{Change of variables}
\begin{theorem}[Change of variables]\label{ch03:thm-cov}
Let $X$ have density $f_X$, let $g$ be differentiable and strictly monotone with inverse $h=g^{-1}$, and
$Y=g(X)$. If $g$ is increasing, $F_Y(y)=F_X(h(y))$ \ts{24}{4}{43:35}; in all cases
\[
  f_Y(y)=f_X\bigl(h(y)\bigr)\,\bigl|h'(y)\bigr| .
\]
\end{theorem}
\begin{proof}
If $g$ is increasing, $\{Y\le y\}=\{X\le h(y)\}$, so $F_Y=F_X\circ h$; differentiating with the chain rule and
$F_X'=f_X$ gives $f_Y=f_X(h)\,h'$ \ts{24}{4}{44:43}. If $g$ is decreasing, $F_Y(y)=\Prob(X\ge h(y))=1-F_X(h(y))$
and $f_Y=-f_X(h)h'=f_X(h)|h'|$.
\end{proof}
The 2013 notes (Theorem~1.3) call $h'$ the \emph{Jacobian} \ts{13}{3}{16:47}; in $\R^n$ it becomes
$f_{\bm Y}(\bm y)=f_{\bm X}(\bm h(\bm y))\,|\det D\bm h(\bm y)|$.

\begin{intuition}[Conservation of probability]
Probability is a conserved ``mass'': $f_Y(y)|\dd y|=f_X(x)|\dd x|$. Where $g$ stretches the axis the density
thins out, where it compresses the axis the density piles up. This is the same bookkeeping as for a
particle density under a change of coordinates.
\end{intuition}

\subsection{The lognormal distribution}
\begin{definition}[Lognormal law]
$Y$ is lognormal$(\mu,\sigma^2)$ if $\ln Y\sim N(\mu,\sigma^2)$, i.e.\ $Y=e^X$ with $X\sim N(\mu,\sigma^2)$
\ts{24}{4}{41:34}. By \cref{ch03:thm-cov} with $h(y)=\ln y$, $h'(y)=1/y$ \ts{24}{4}{45:45}:
\[
  f_Y(y)=\frac{1}{y\,\sigma\sqrt{2\pi}}\exp\Bigl(-\frac{(\ln y-\mu)^2}{2\sigma^2}\Bigr),\qquad y>0
\]
\ts{24}{4}{46:56}\ts{13}{3}{21:12}.
\end{definition}

\noindent Properties (the first three from monotonicity of $\ln$ and a derivative):
\begin{itemize}
  \item quantiles map to quantiles: $q_p(Y)=e^{\mu+\sigma z_p}$ with $z_p=\Phi^{-1}(p)$; the median is $e^\mu$;
  \item the mode is $e^{\mu-\sigma^2}$;
  \item $\E[Y^k]=e^{k\mu+k^2\sigma^2/2}$ for every real $k$ (\cref{ch03:sec-mgf}); in particular
        $\E [Y]=e^{\mu+\sigma^2/2}$, so mode $<$ median $<$ mean: the law is skewed to the right;
  \item products and powers: if $Y_1,Y_2$ are independent lognormal, $Y_1Y_2$ and $Y_1^c$ are lognormal.
\end{itemize}
The parameters $\mu,\sigma$ are the mean and standard deviation of $\ln Y$, \emph{not} of $Y$
\ts{13}{3}{23:15}; they are kept because the law is defined through the logarithm \ts{13}{3}{24:18}.

\begin{example}[Lognormal$(0.2,0.4^2)$, slides 15--16 of \file{lec04_1}]\label{ch03:ex-ln}
Let $X_0=1$ and $X_1$ lognormal with $\mu=0.2$, $\sigma=0.4$: a one-year log-return with mean $0.2$ and
standard deviation $0.4$ \ts{24}{4}{47:56}. Then (\cref{ch03:fig-lognormal}) the 5\%, 25\%, 50\%, 75\%, 95\%
quantiles are $0.633$, $0.933$, $1.221$, $1.600$, $2.358$ \ts{24}{4}{48:58}; the mode is $1.041$, the mean
$1.323$ and the standard deviation $0.551$. The quantiles are symmetric \emph{multiplicatively} about the
median: $1.221/0.633=2.358/1.221=1.93$.
\end{example}

\begin{figure}[ht]
\centering
\begin{tikzpicture}
\begin{axis}[width=0.85\linewidth,height=5.6cm,axis lines=left,xmin=0,xmax=3,ymin=0,ymax=1.12,
  xlabel={$y$},ylabel={$f_Y(y)$},tick label style={font=\small},label style={font=\small}]
  \addplot[very thick,defcol,domain=0.03:3,samples=150]
    {exp(-(ln(x)-0.2)^2/(2*0.16))/(x*0.4*sqrt(2*pi))};
  \draw[thmcol,thick] (axis cs:0.6326,0) -- (axis cs:0.6326,0.92) node[above,font=\scriptsize]{5\%};
  \draw[intcol,thick] (axis cs:0.9326,0) -- (axis cs:0.9326,0.92) node[above,font=\scriptsize]{25\%};
  \draw[black,thick] (axis cs:1.2214,0) -- (axis cs:1.2214,1.0) node[above,font=\scriptsize]{median};
  \draw[intcol,thick] (axis cs:1.5997,0) -- (axis cs:1.5997,0.92) node[above,font=\scriptsize]{75\%};
  \draw[thmcol,thick] (axis cs:2.3583,0) -- (axis cs:2.3583,0.92) node[above,font=\scriptsize]{95\%};
  \draw[fincol,thick,dashed] (axis cs:1.3231,0) -- (axis cs:1.3231,0.62);
  \node[font=\scriptsize,fincol,anchor=west] at (axis cs:1.33,0.06) {mean};
\end{axis}
\end{tikzpicture}
\caption{Density of the lognormal$(0.2,0.4^2)$ law with its 5, 25, 50, 75, 95\% quantiles (as on slide~16
of \file{lec04_1}) and its mean. The mode ($1.041$, the peak) lies left of the median, the mean right of
it.}\label{ch03:fig-lognormal}
\end{figure}

\begin{coursenote}[``An annualised return of about 20\%'']
In class the example is read as an asset with an annualised return of about 20\% and volatility 0.4
\ts{24}{4}{47:56}. Be careful which return: $0.2$ is the expected \emph{log}-return; the median gross
return is $e^{0.2}=1.221$ (+22.1\%) and the expected gross return is $e^{0.28}=1.323$ (+32.3\%). The gap
$\sigma^2/2=0.08$ is the convexity correction that reappears as the $-\sigma^2/2$ in the solution of the
geometric Brownian motion (\cref{ch:stochcalc}).
\end{coursenote}

\begin{finance}[Log-returns and the square-root-of-time rule]
Log-returns add over time: $\ln(S_T/S_0)=\sum_{\text{days}}r_i$. If daily log-returns are i.i.d.\ with mean
$\mu_d$ and variance $\sigma_d^2$, the $T$-day log-return has mean $T\mu_d$ and standard deviation
$\sigma_d\sqrt T$. This is why volatilities are quoted annualised and scaled by $\sqrt t$ (as in 2024 PS2
Problem~2(c), where $\sigma=\sqrt t\,\ln1.3$), and why the Black--Scholes model takes $S_T$ lognormal
(\cref{ch:bs}). For small moves the simple return $e^{r}-1\approx r$.
\end{finance}\begin{answer}{vorrei che la notazione fosse consistente per tutto il libro. S e' il prezzo della stock, giusto? cos'e' P allora? B e' il prezzo del bond?}[$S$, $P$, $B$]
Yes: $S$ = stock price, $B$ = bond (risk-free asset) price. $P$ is a \emph{generic} asset price, any asset;
above it comes from the 2013 lecture, which uses $P_n$ for the price on day $n$ of an unspecified asset
(as $P^j$ in \cref{def:onepmodel}). Here $P_n$ and $S_n$ are the same thing.
\end{answer}

\begin{exercise}[PS2 (2013), B-2]\label{ch03:ex-b2}
\leavevmode
\begin{enumerate}[(a)]
  \item Verify that $\phi(x)=\frac1{\sigma\sqrt{2\pi}}e^{-(x-\mu)^2/(2\sigma^2)}$ is a probability density (compute
      $\bigl(\int\phi\bigr)^2$ in polar coordinates).
  \item Compute the mean and variance of the lognormal law (it suffices to
      treat $\mu=0$; use (a) to simplify). (This is also Exercise~1.5(i) of \file{lecnote3}; Exercise~1.2 there asks
      to verify that the mean of $N(\mu,\sigma^2)$ is $\mu$.)
\end{enumerate}
\end{exercise}

\begin{exercise}[PS2 (2013), B-4: does time diversify risk?]
(From M.~Kritzman, \emph{Puzzles of Finance}, ch.~3; see also Kritzman and Rich, \emph{Beware of dogma: the truth
about time diversification}.) A stock has independent annual returns, each lognormal with $\mu=\ln1.1$ and
$\sigma=\ln1.2$. According to the author, the fraction of wealth lost with 0.1\% probability over $T$ years is
40.5\%, 58.43\%, 61.48\% and 73.92\% for $T=1,5,10,20$, and he concludes that if risk is measured by the
magnitude of potential loss, it \emph{increases} with the horizon.
\begin{enumerate}[(a)]
  \item Compute the fraction $h(T)$ of wealth
      lost with probability 0.1\% over $T$ years (in terms of $\Phi$).
  \item Find the $T$ that maximises $h(T)$. Is $h$
      increasing in $T$?
  \item Criticise the argument using (a) and (b).
\end{enumerate}
\end{exercise}

\begin{exercise}[\file{lec04_1}, slide 21]
\leavevmode
\begin{enumerate}[(a)]
  \item Find the skewness of a lognormal$(\mu,\sigma^2)$ variable (use the lognormal moments).
  \item Show directly
      from the density that the kurtosis of $N(\mu,\sigma^2)$ is 3.
\end{enumerate}
\end{exercise}

% ------------------------------------------------------------------
\section{Expected option payoffs}\label{ch03:sec-call}

A European call gives the right to buy the asset at time $T$ at the strike $K$; if $X$ is the price at $T$,
the payoff is the ``hockey stick'' $C=(X-K)^+=\max(0,X-K)$ \ts{24}{4}{50:01} (\cref{ch03:fig-hockey}).\note*{plot it}

\begin{figure}[!htb]
\centering
\begin{tikzpicture}
\begin{axis}[width=6.6cm,height=4.6cm,axis lines=middle,xmin=0,xmax=2.05,ymin=-0.1,ymax=1.1,
  xtick={1},xticklabels={$K$},ytick=\empty,xlabel={$X$},tick label style={font=\scriptsize},label style={font=\small},title style={font=\small},legend style={font=\scriptsize},title={call: $(X-K)^+$},
  xlabel style={at={(1,0)},anchor=north}]
  \addplot[very thick,intcol] coordinates {(0,0) (1,0) (2,1)};
\end{axis}
\end{tikzpicture}\hfill
\begin{tikzpicture}
\begin{axis}[width=6.6cm,height=4.6cm,axis lines=middle,xmin=0,xmax=2.05,ymin=-0.1,ymax=1.1,
  xtick={1},xticklabels={$K$},ytick={1},yticklabels={$K$},xlabel={$X$},tick label style={font=\scriptsize},label style={font=\small},title style={font=\small},legend style={font=\scriptsize},title={put: $(K-X)^+$},
  xlabel style={at={(1,0)},anchor=north}]
  \addplot[very thick,thmcol] coordinates {(0,1) (1,0) (2,0)};
\end{axis}
\end{tikzpicture}
\caption{Payoffs at expiry as functions of the price $X$ of the underlying: the call (the ``hockey stick'': flat
blade, then slope 1) and the put, which pays at most $K$.}\label{ch03:fig-hockey}
\end{figure}

\begin{theorem}[Expected call payoff]\label{ch03:thm-call}
Let $X$ have density $f$, CDF $F$ and $\E[|X|]<\infty$. For every constant $K$,
\[
  \E\bigl[(X-K)^+\bigr]=\int_K^\infty\!\!\int_x^\infty f(t)\dd t\dd x=\int_K^\infty\bigl[1-F(x)\bigr]\dd x .
\]
\end{theorem}
\begin{proof}
Following the slides \ts{24}{4}{51:05}, for $M>K$ integrate by parts with $u=x-K$, $v=-[1-F(x)]$:
\begin{align*}
  \int_K^M(x-K)f(x)\dd x&=-(M-K)\bigl[1-F(M)\bigr]\\
  &\quad+\int_K^M\bigl[1-F(x)\bigr]\dd x .
\end{align*}
The boundary term vanishes as $M\to\infty$ because $0\le(M-K)[1-F(M)]\le\int_M^\infty(x-K)f(x)\dd x$, the tail
of a convergent integral. (Alternatively: $(x-K)^+=\int_K^\infty\ind\{x>s\}\dd s$; take expectations and swap
the integrals.)
\end{proof}
This is the familiar fact that for a non-negative variable $\E[W]=\int_0^\infty\Prob(W>s)\dd s$
\ts{24}{4}{52:11}, applied to $W=(X-K)^+$. The slides assume finite variance; a finite mean is enough.

\begin{corollary}[Strike derivatives; not discussed in class]\label{ch03:cor-bl}
$C(K)=\E[(X-K)^+]$ satisfies $C'(K)=-\Prob(X>K)$ and $C''(K)=f(K)$.
\end{corollary}
So call prices across strikes determine the distribution of $X$ (under the pricing measure): this is
the Breeden--Litzenberger formula behind \cref{ch:duality,ch:ross}.\begin{answer}{questa parte non l'ho capita}[Breeden--Litzenberger]
Read $C(K)$ as the (undiscounted) price of a call, for every strike $K$. Differentiate twice in $K$: the
corollary gives $C''(K)=f(K)$. So if the market quotes calls at many strikes, the curvature of the price
curve \emph{is} the density of $S_T$ (under $Q$, \cref{ans:pricingQ}). Intuition: a ``butterfly''
(buy calls at $K-h$ and $K+h$, sell two at $K$) pays a small triangle around $K$, i.e.\ roughly
$h$ dollars if $S_T\approx K$; its price $\approx h^2C''(K)$ is the price of ``$S_T$ near $K$''.
\end{answer}

\begin{corollary}[Normal and lognormal payoffs, slide 18 of \file{lec04_1}]\label{ch03:cor-normal}
\leavevmode
\begin{enumerate}[(i)]
  \item If $X\sim N(\mu,\sigma^2)$ and $d=(\mu-K)/\sigma$, then
      \[
      \E\bigl[(X-K)^+\bigr]=(\mu-K)\,\Phi(d)+\sigma\,\phi(d)=(\mu-K)\,\Phi(d)+\frac{\sigma}{\sqrt{2\pi}}\,e^{-(K-\mu)^2/(2\sigma^2)} .
      \]
  \item If $X$ is lognormal$(\mu,\sigma^2)$ and $d=(\mu-\ln K)/\sigma$, then
      \[
      \E\bigl[(X-K)^+\bigr]=e^{\mu+\sigma^2/2}\,\Phi(d+\sigma)-K\,\Phi(d)
      \]
      \ts{24}{4}{53:14}.
\end{enumerate}
\end{corollary}\begin{answer}{non c'e' un errore nell'ultima uguaglianza di (i)?}[No error]
It is correct: $\phi(d)=e^{-d^2/2}/\sqrt{2\pi}$ with $d^2=(\mu-K)^2/\sigma^2$, and $(\mu-K)^2=(K-\mu)^2$.
\end{answer}
\begin{proof}
(i) Write $X=\mu+\sigma Z$: $(X-K)^+=(\mu-K+\sigma Z)\ind\{Z>-d\}$. Then $\Prob(Z>-d)=\Phi(d)$ and
$\E[Z\ind\{Z>a\}]=\int_a^\infty z\phi(z)\dd z=\phi(a)$, with $\phi(-d)=\phi(d)$.

(ii) is 2024 PS2 Problem~2(d)
(\cref{ch03:ex-ps2}).
\end{proof}

\begin{finance}[From expected payoffs to option prices]
A price is a \emph{discounted} expectation under a \emph{pricing} measure (\cref{thm:ftap}):\begin{answer}{cos'e' una pricing measure Q?}[Pricing measure]\label{ans:pricingQ}
The weights $q_i=\psi_i/\sum\psi$ of \cref{def:onepmodel}: positive numbers summing to 1, so they
\emph{look} like probabilities, chosen so that every price is a discounted average of payoffs,
$P_0=\beta\E_Q[P_T]$. They are not the real chances: they also include risk aversion (bad states get more
weight). Also called the \emph{risk-neutral} measure, see \cref{ch:bs}.
\end{answer}
$C_0=e^{-rT}\E_Q[(S_T-K)^+]$. In the Black--Scholes model, under $Q$, $\ln S_T\sim N(\ln S_0+(r-\sigma^2/2)T,\
\sigma^2T)$; inserting these parameters in \cref{ch03:cor-normal}(ii) gives $e^{\mu+\sigma^2T/2}=S_0e^{rT}$, $d=d_2$,
$d+\sigma\sqrt T=d_1$ and hence $C_0=S_0\Phi(d_1)-Ke^{-rT}\Phi(d_2)$ (\cref{ch:bs}). Part~(1) is the
\emph{Bachelier} model (1900), still used for interest rates and spreads, which can be negative. At the
money ($K=\mu$) it gives $\sigma/\sqrt{2\pi}\approx0.4\,\sigma$: the traders' rule of thumb ``an at-the-money
option is worth about 0.4 standard deviations of the underlying at expiry'' (not discussed in class).\begin{answer}{can you give a brief example?}[Rule of thumb]
A stock at \$100 with volatility 20\%/year; a 3-month at-the-money option. The standard deviation of $S_T$
is about $100\cdot0.2\cdot\sqrt{0.25}=\$10$, so the option is worth about $0.4\cdot10=\$4$. Black--Scholes
($r=0$) gives \$3.99.
\end{answer}
\end{finance}

\begin{exercise}[not from the course: puts and parity]
With the notation of \cref{ch03:thm-call}, show that $\E[(K-X)^+]=\int_{-\infty}^KF(x)\dd x$ and that
$\E[(X-K)^+]-\E[(K-X)^+]=\E[X]-K$. Interpret the second identity as put--call parity: a long call plus a short put is a
forward contract (\cref{sec:forward}).
\end{exercise}

% ------------------------------------------------------------------
\section{Moment-generating functions}\label{ch03:sec-mgf}

\begin{definition}[MGF, characteristic function]
The \emph{moment-generating function} of $X$ is $M_X(t)=\E[e^{tX}]$, $t\in\R$ \ts{24}{4}{53:14}\ts{13}{3}{36:02}.
It always satisfies $M_X(0)=1$, and the set where it is finite is an interval containing $0$ (by
convexity of $t\mapsto e^{tx}$). We say that $X$ \emph{has an MGF} if $M_X$ is finite on some
$(-\delta,\delta)$. The \emph{characteristic function} $\varphi_X(t)=\E[e^{itX}]=\E[\cos tX]+i\,\E[\sin tX]$
is finite for every $X$ and $t$, since $|e^{itX}|=1$ \ts{24}{4}{56:28}.
\end{definition}\begin{answer}{they are related by analitical continuation. posso dire che una e' la laplace transform e l'altra la fourier transform?}[Yes]
Yes: $M_X(t)$ is the (two-sided) Laplace transform of the density at $-t$, $\varphi_X(t)$ its Fourier
transform (up to the sign convention), and $\varphi_X(t)=M_X(it)$ when $M_X$ is finite near 0. As in
physics, the Fourier version always exists, the Laplace one only with enough decay.
\end{answer}

\begin{proposition}[Moments from the MGF]
If $M_X$ is finite on $(-\delta,\delta)$, all moments are finite, $\E[X^k]=M_X^{(k)}(0)$ and
$M_X(t)=\sum_{k\ge0}\frac{t^k}{k!}\,m_k$ for $|t|<\delta$ \ts{13}{3}{38:14}\ts{24}{4}{54:18}.
\end{proposition}
\begin{proof}
For $|s|<\delta$, $e^{s|x|}\le e^{sx}+e^{-sx}$, so $\E[e^{s|X|}]<\infty$; this dominates $\sum_k|tX|^k/k!$ for
$|t|\le s$, so expectation and series (and derivatives) can be exchanged.
\end{proof}

\paragraph{The MGF need not exist} \ts{24}{4}{55:21}. For the Cauchy density $f(x)=1/(\pi(1+x^2))$ (\cref{ch03:fig-laws}),\note*{plot it} $\E[|X|]=\infty$
and $M_X(t)=\infty$ for every $t\ne0$; its characteristic function is $e^{-|t|}$. The \emph{lognormal} law has
all moments but no MGF \ts{13}{3}{37:09}: for $t>0$ and every $k$, $e^{tY}\ge(tY)^k/k!$, hence
\[
  M_Y(t)\ \ge\ \frac{t^k\,\E[Y^k]}{k!}=\frac{t^k e^{k\mu+k^2\sigma^2/2}}{k!}\ \ge\ \exp\Bigl(\frac{k^2\sigma^2}{2}+k(\mu+\ln t)-k\ln k\Bigr)\xrightarrow[k\to\infty]{}\infty ,
\]
using $k!\le k^k$. For $t\le0$, instead, $0<e^{tY}\le1$ and $M_Y(t)$ is finite.

\begin{coursenote}[Erratum in the 2013 notes]
\file{lecnote3} (and the lecture, \ts{13}{3}{46:47}) state that for the lognormal law $\E[e^{tX}]$ does not
exist for all $t\neq0$. It is infinite for all $t>0$ but finite for all $t<0$; what fails is finiteness on a
neighbourhood of $0$, which is what the theorems below need. In the same paragraph,
$\frac{d^kM_X}{dx^k}(0)$ should read $\frac{d^kM_X}{dt^k}(0)$.
\end{coursenote}

\begin{theorem}[Uniqueness and continuity]\label{ch03:thm-mgf}
\leavevmode
\begin{enumerate}[(i)]
  \item If $M_X(t)=M_Y(t)<\infty$ for all $t$ in a neighbourhood of $0$, then $X$ and $Y$ have the same distribution
        \ts{24}{4}{56:28}\ts{13}{3}{39:18}.
  \item If $M_{X_n}(t)\to M_X(t)<\infty$ for all $t$ in a neighbourhood of $0$, then $X_n$ converges to $X$ \emph{in
        distribution}: $F_{X_n}(x)\to F_X(x)$ at every point where $F_X$ is continuous \ts{24}{4}{57:31}\ts{13}{3}{43:39}.
\end{enumerate}
The same statements hold for characteristic functions, with ``for all $t\in\R$'' and no existence
assumption. (Not proved in class; see Durrett, \emph{Probability: Theory and Examples}.)
\end{theorem}

\begin{coursenote}[Statement in the slides and notes]
Slide~19 of \file{lec04_1} and Theorem~2.2 of \file{lecnote3} assume equality or convergence ``for all $t$''.
The neighbourhood version above is the standard one and is the one needed in practice: the MGF of a
$\chi^2$ variable, for instance, is finite only for $t<1/2$ (\cref{ch03:sec-chisq}).
\end{coursenote}\begin{answer}{per ovviare al fatto che alcune distribuzioni non hanno i momenti ben definiti, perche' non aggiungiamo un weight esponenziale all'integrale cosi' abbiamo quantita' finite per caratterizzare le nostre distribuzioni heavy-tailed?}[It is done]
Good idea, and it is used. A damping $e^{-a|x|}$ makes every ``moment'' finite, but the numbers then
depend on $a$ and lose their meaning (mean, variance). The bounded weight $e^{itx}$ is the standard
answer: the characteristic function always exists and determines the law. Practitioners also use
quantile-based measures (median, interquartile range) and the tail index $\alpha$ of $\Prob(|X|>x)\sim x^{-\alpha}$.
In option pricing, damped transforms $e^{-ax}$ are the Carr--Madan (1999) method of pricing options by FFT.
\end{answer}

\begin{remark}[Moments do not determine a distribution]
If two variables have the same moments of all orders, they need not have the same law: the hole in the
argument ``same moments $\Rightarrow$ same MGF $\Rightarrow$ same law'' is that the MGF may not exist
\ts{13}{3}{41:26}\ts{13}{3}{42:28}. The standard example (Durrett, pp.~106--107, cited by \file{lecnote3}) is the
lognormal. Let $f$ be the lognormal$(0,1)$ density and, for $|a|\le1$, $f_a(y)=f(y)\,[1+a\sin(2\pi\ln y)]\ge0$.
With $y=e^u$, $\phi(u)e^{ku}=e^{k^2/2}\phi(u-k)$ and $v=u-k$,
\[
  \int_0^\infty y^kf(y)\sin(2\pi\ln y)\dd y=e^{k^2/2}\int_{\R}\phi(v)\sin(2\pi v+2\pi k)\dd v=0
  \qquad(k=0,1,2,\dots),
\]
since $\sin$ is odd. So every $f_a$ is a density with exactly the lognormal moments. (Not discussed in class.)
\end{remark}\begin{answer}{pero' non c'e' il problema che la pdf non e' definita positiva? inoltre, i momenti non si possono ricavare dalla characteristic function che esiste sempre?}[Positivity]
No: $|a\sin(\cdot)|\le1$, so $1+a\sin(2\pi\ln y)\ge0$ and $f_a\ge0$; it integrates to 1 (the $k=0$ case).
\answerpart{Moments from $\varphi$?}
Yes, $\E[X^k]=i^{-k}\varphi^{(k)}(0)$ when the moments exist. But the converse fails: the moments give only
the Taylor coefficients of $\varphi$ at 0, and for the lognormal that Taylor series has radius 0, so it
does not reconstruct $\varphi$. Different laws, different $\varphi$, same derivatives at 0.
\end{answer}

\paragraph{The normal MGF.} For $Z\sim N(0,1)$, complete the square, $tx-\tfrac12x^2=-\tfrac12(x-t)^2+\tfrac12t^2$
\ts{24}{4}{58:34}\ts{24}{4}{1:00:40}:
\[
  M_Z(t)=\int_{\R}\frac{e^{-\frac12(x-t)^2}}{\sqrt{2\pi}}\dd x\;e^{t^2/2}=e^{t^2/2}.
\]
Since $M_{a+bX}(t)=e^{at}M_X(bt)$ \ts{24}{4}{1:01:42}, for $X=\mu+\sigma Z\sim N(\mu,\sigma^2)$:
\begin{equation}\label{ch03:eq-normalmgf}
  M_X(t)=e^{\mu t+\sigma^2t^2/2}.
\end{equation}
Consequences:
\begin{itemize}
  \item \emph{Lognormal moments} \ts{24}{4}{1:02:42}: $\E[e^{kX}]=M_X(k)$, so $\E[Y^k]=e^{k\mu+k^2\sigma^2/2}$ for
        $Y=e^X$ (2024 PS2 Problem~2(a,b)).
  \item \emph{Sums}: if $X_1,\dots,X_n$ are independent, $M_{\sum X_i}(t)=\prod_iM_{X_i}(t)$, because the expectation of a
        product of independent variables is the product of expectations. Hence a sum of independent normals
        is $N(\sum\mu_i,\sum\sigma_i^2)$, by \cref{ch03:thm-mgf}(i).
  \item \emph{Normal moments}: $e^{t^2/2}=\sum_k t^{2k}/(2^kk!)$ gives $\E[Z^{2k}]=(2k)!/(2^kk!)$, e.g.\ $\E[Z^4]=3$:
        the normal kurtosis is $3$.
\end{itemize}

\begin{intuition}[Partition function and cumulants (not discussed in class)]
$M_X(t)$ is a Laplace transform, like a partition function $Z(\beta)$, and $K_X(t)=\ln M_X(t)$ plays the role of
a free energy. Its Taylor coefficients are the \emph{cumulants}: $K_X(t)=\mu t+\sigma^2\frac{t^2}{2}+\kappa_3\frac{t^3}{6}+
\kappa_4\frac{t^4}{24}+\cdots$ with $\kappa_3=\gamma\sigma^3$ and $\kappa_4=(\kappa-3)\sigma^4$. For independent variables
cumulants \emph{add}. The Gaussian is the ``free theory'': $K$ is exactly quadratic. For the standardised sum
$S_n/\sqrt n$ of i.i.d.\ variables the $k$-th cumulant is $n\kappa_k/n^{k/2}\to0$ for $k\ge3$: all non-Gaussian
features wash out, which is the central limit theorem in one line (\cref{ch03:thm-clt}).
\end{intuition}\begin{answer}{beautiful box. i want to know more.}[More on cumulants]
Three directions:
\begin{itemize}
  \item \emph{Large deviations} (Cram\'er): $\Prob(S_n/n\approx x)\approx e^{-nI(x)}$ with $I$ the Legendre
        transform of $K$, exactly free energy $\leftrightarrow$ entropy.
  \item \emph{Edgeworth expansion}: corrections to the CLT in powers of $\kappa_3/\sqrt n$, $\kappa_4/n$.
  \item \emph{Exponential tilting} $f_\theta(x)=e^{\theta x-K(\theta)}f(x)$ (the canonical ensemble): the Esscher
        transform, a way to build a pricing measure.
\end{itemize}
Reference: Feller, vol.~2, ch.~XVI; Dembo--Zeitouni, \emph{Large Deviations}.
\end{answer}

\paragraph{Exponential families (2013).} A family of pdfs or pmfs indexed by a parameter vector $\theta$ is an
\emph{exponential family} if \ts{13}{3}{26:24}
\[
  f(x\mid\theta)=h(x)\,c(\theta)\exp\Bigl(\sum_{i=1}^kw_i(\theta)\,t_i(x)\Bigr),
\]
with $h,t_i$ depending only on $x$ and $c,w_i$ only on $\theta$. The lognormal is one \ts{13}{3}{29:37}:
expanding the square,
\[
  f(x)=\frac1x\cdot\frac{e^{-\mu^2/(2\sigma^2)}}{\sigma\sqrt{2\pi}}\cdot\exp\Bigl(-\frac{(\ln x)^2}{2\sigma^2}+\frac{\mu\ln x}{\sigma^2}\Bigr),
\]
so $h(x)=1/x$, $c(\theta)=e^{-\mu^2/(2\sigma^2)}/(\sigma\sqrt{2\pi})$, $w_1=-1/(2\sigma^2)$, $t_1=(\ln x)^2$, $w_2=\mu/\sigma^2$,
$t_2=\ln x$. The point, stressed by P.~Kempthorne from the audience \ts{13}{3}{28:35}: the joint density of an
i.i.d.\ sample $\prod_jf(x_j\mid\theta)=\prod_jh(x_j)\,c(\theta)^n\exp\bigl(\sum_iw_i(\theta)\sum_jt_i(x_j)\bigr)$ has the
same form and depends on the data only through the \emph{sufficient statistics} $\sum_jt_i(x_j)$. Normal,
lognormal, exponential, Poisson and gamma laws are exponential families (2013 PS2 B-3); Cauchy and Student $t$
are not.

\begin{coursenote}[Erratum in the 2013 notes]
In the lognormal example of \file{lecnote3} (p.~3) the factor is printed $h(x)=\frac12$; it should be
$h(x)=\frac1x$, as in the displayed formula just above it.
\end{coursenote}

\begin{table}[ht]
\centering\small
\begin{tabularx}{\linewidth}{@{}l l L@{}}
\toprule
\textbf{Law} & \textbf{pmf / pdf} & \textbf{Where it appears}\\
\midrule
Bernoulli$(p)$ & $p^x(1-p)^{1-x}$, $x\in\{0,1\}$ & default indicators, side of an order, event contracts (\cref{ch:kalshi})\\
Poisson$(\lambda)$ & $\lambda^ke^{-\lambda}/k!$, $k\in\N$ & number of events (orders, jumps, defaults) in a period; 2013 PS2 A-2, B-1\\
Exponential$(\lambda)$ & $\lambda e^{-\lambda x}$, $x\ge0$ & waiting times between events (memoryless); 2013 PS2 A-1\\
Normal$(\mu,\sigma^2)$ & $\frac{1}{\sigma\sqrt{2\pi}}e^{-(x-\mu)^2/(2\sigma^2)}$ & log-returns, Brownian increments, regression errors\\
Lognormal$(\mu,\sigma^2)$ & $\frac{1}{x\sigma\sqrt{2\pi}}e^{-(\ln x-\mu)^2/(2\sigma^2)}$, $x>0$ & prices (Black--Scholes)\\
Gamma$(\alpha,\lambda)$ & $\frac{\lambda^\alpha}{\Gamma(\alpha)}x^{\alpha-1}e^{-\lambda x}$, $x>0$ & contains $\chi^2_n=\mathrm{Gamma}(n/2,1/2)$\\
$\chi^2_n$, $t_r$, $F_{m,n}$ & \cref{ch03:sec-chisq} & inference in regression (\cref{ch:regression}); $t$: heavy-tailed returns\\
Cauchy & $1/(\pi(1+x^2))$ & a law without mean, variance or MGF\\
\bottomrule
\end{tabularx}
\caption{Distributions used in the course.}\label{ch03:tab-laws}
\end{table}\note*{plot all the mentioned distributions}

\begin{figure}[!htb]
\centering
\pgfplotsset{lawplot/.style={width=5.3cm,height=4.2cm,tick label style={font=\scriptsize},label style={font=\small},title style={font=\small},legend style={font=\scriptsize},legend pos=north east}}
\begin{tikzpicture}
\begin{axis}[lawplot,title={Bernoulli$(0.3)$, Poisson$(3)$},ybar=0pt,bar width=4pt,xmin=-0.7,xmax=9.7,ymin=0]
  \addplot[fill=defcol!60,draw=none] coordinates {(0,0.7) (1,0.3)};
  \addplot[fill=thmcol!60,draw=none] coordinates {(0,0.0498) (1,0.1494) (2,0.2240) (3,0.2240) (4,0.1680) (5,0.1008) (6,0.0504) (7,0.0216) (8,0.0081) (9,0.0027)};
  \legend{Bernoulli,Poisson}
\end{axis}
\end{tikzpicture}\hfill
\begin{tikzpicture}
\begin{axis}[lawplot,title={exponential, gamma, $\chi^2$},xmin=0,xmax=8,ymin=0,ymax=1.05]
  \addplot[thick,defcol,domain=0:8,samples=100] {exp(-x)};
  \addplot[thick,intcol,domain=0.01:8,samples=100] {x*exp(-x)};
  \addplot[thick,thmcol,domain=0.05:8,samples=100] {x^0.5*exp(-x/2)/(2^1.5*0.886227)};
  \legend{Exp$(1)$,{Gamma$(2,1)$},$\chi^2_3$}
\end{axis}
\end{tikzpicture}\hfill
\begin{tikzpicture}
\begin{axis}[lawplot,title={normal, $t_3$, Cauchy},xmin=-5,xmax=5,ymin=0,ymax=0.42]
  \addplot[thick,defcol,domain=-5:5,samples=120] {exp(-x^2/2)/sqrt(2*pi)};
  \addplot[thick,intcol,domain=-5:5,samples=120] {0.367553*(1+x^2/3)^(-2)};
  \addplot[thick,thmcol,domain=-5:5,samples=120] {1/(pi*(1+x^2))};
  \legend{{$N(0,1)$},$t_3$,Cauchy}
\end{axis}
\end{tikzpicture}
\begin{tikzpicture}
\begin{axis}[lawplot,title={lognormal$(0,\sigma^2)$},xmin=0,xmax=4,ymin=0,ymax=1.7]
  \addplot[thick,defcol,domain=0.02:4,samples=120] {exp(-ln(x)^2/(2*0.0625))/(x*0.25*sqrt(2*pi))};
  \addplot[thick,intcol,domain=0.02:4,samples=120] {exp(-ln(x)^2/(2*0.25))/(x*0.5*sqrt(2*pi))};
  \addplot[thick,thmcol,domain=0.02:4,samples=120] {exp(-ln(x)^2/2)/(x*sqrt(2*pi))};
  \legend{$\sigma=0.25$,$\sigma=0.5$,$\sigma=1$}
\end{axis}
\end{tikzpicture}\hfill
\begin{tikzpicture}
\begin{axis}[lawplot,title={tails, log scale},xmin=0,xmax=8,ymode=log,ymin=1e-6,ymax=1,legend pos=south west]
  \addplot[thick,defcol,domain=0:5.3,samples=80] {exp(-x^2/2)/sqrt(2*pi)};
  \addplot[thick,intcol,domain=0:8,samples=80] {0.367553*(1+x^2/3)^(-2)};
  \addplot[thick,thmcol,domain=0:8,samples=80] {1/(pi*(1+x^2))};
  \legend{{$N(0,1)$},$t_3$,Cauchy}
\end{axis}
\end{tikzpicture}\hfill
\begin{tikzpicture}
\begin{axis}[lawplot,title={$F_{m,n}$},xmin=0,xmax=4,ymin=0,ymax=1]
  \addplot[thick,defcol,domain=0.01:4,samples=120] {10.368*x^1.5*(1+0.5*x)^(-7.5)};
  \legend{$F_{5,10}$}
\end{axis}
\end{tikzpicture}
\caption{The laws of \cref{ch03:tab-laws}. The $t_3$ and Cauchy densities look like the normal near $0$ but
their tails decay like powers ($x^{-4}$, $x^{-2}$), not like $e^{-x^2/2}$: on the log scale the gap at $x=5$ is a
factor $10^3$--$10^4$.}\label{ch03:fig-laws}
\end{figure}

\begin{exercise}[PS2 (2024), Problem 2: lognormal distribution]\label{ch03:ex-ps2}
Let $X$ be lognormal$(\mu,\sigma^2)$, i.e.\ $Y=\ln X\sim N(\mu,\sigma^2)$.
\begin{enumerate}[(a)]
  \item Prove that $\mu_*=\E[X]=e^{\mu+\sigma^2/2}$.
  \item Prove that $\sigma_*^2=\Var X=e^{2\mu+2\sigma^2}-e^{2\mu+\sigma^2}=\mu_*^2\,(e^{\sigma^2}-1)$.
      For (a) and (b) use the MGF of the normal law; no integrals are needed.
  \item A stock has price $S_0$ at $t=0$ and $S_0X$ at time $t$ (years), where $X$ is lognormal with $\mu=t\ln1.15$
      and $\sigma=\sqrt t\,\ln1.3$. Find the mean and standard deviation of the total return $X$ for $t=1/12,\,3/12,\,1$.
  \item Show that for a call with strike $K$ on a lognormal$(\mu,\sigma^2)$ price,
      $\E[(X-K)^+]=e^{\mu+\sigma^2/2}\Phi\bigl(\frac{\mu-\ln K}{\sigma}+\sigma\bigr)-K\Phi\bigl(\frac{\mu-\ln K}{\sigma}\bigr)$.
  \item For the stock of (c) with $S_0=\$100$, compute the expected payoff of the call for $T=1/12,\,3/12,\,1$ and
      strikes $K=\$100$ (at the money) and $K=\$120$ (out of the money). Comment on the effect of a longer time to
      expiration.
\end{enumerate}
\end{exercise}

\begin{exercise}[PS2 (2013), A-1: exponential distribution]
$X$ has density $f_X(x)=\lambda e^{-\lambda x}$, $x\ge0$.
\begin{enumerate}[(a)]
  \item Compute its MGF.
  \item Deduce mean and variance from the MGF.
  \item Compute $\Prob(X>t)$ for $t>0$.
  \item Prove the memoryless property $\Prob(X>s+t\mid X>s)=\Prob(X>t)$ for all $s,t>0$.
  \item If $X_1,\dots,X_n$ are i.i.d.\ exponential$(\lambda)$, prove that $\min\{X_1,\dots,X_n\}$ is exponential.
  \item Entering a bank you find all three tellers busy and nobody waiting. Service times of the three customers and
      yours are i.i.d.\ exponential. What is the probability that you are the last of the four to leave?
\end{enumerate}
\end{exercise}

\begin{exercise}[PS2 (2013), A-2: Poisson distribution]
$X$ has pmf $f_X(k)=\lambda^ke^{-\lambda}/k!$, $k=0,1,2,\dots$
\begin{enumerate}[(a)]
  \item Compute its MGF.
  \item Deduce mean and variance.
  \item Prove, with MGFs, that a sum of independent Poisson
      variables is Poisson.
\end{enumerate}
\end{exercise}

\begin{exercise}[PS2 (2013), B-1: from exponential waiting times to Poisson counts]
Let $X_1,X_2,\dots$ be i.i.d.\ exponential$(\lambda)$ and $\tau=\max\{t\ge0:X_1+\dots+X_t\le1\}$, the number of events with
exponential waiting times in an interval of length 1.
\begin{enumerate}[(a)]
  \item Compute $\Prob(\tau=0)$.
  \item Compute $\Prob(\tau=n)$ using
      that $S_n=X_1+\dots+X_n$ has CDF $F_{S_n}(x)=1-\sum_{k=0}^{n-1}\frac1{k!}e^{-\lambda x}(\lambda x)^k$ ($x>0$), and conclude that
      $\tau$ is Poisson.
\end{enumerate}
\end{exercise}

\begin{exercise}[PS2 (2013), B-3: exponential families]
Prove that the lognormal, Poisson and exponential distributions are exponential families, i.e.\ that their
densities or pmfs can be written as $h(x)c(\theta)\exp\bigl(\sum_iw_i(\theta)t_i(x)\bigr)$ with $c\ge0$, $h\ge0$.
\end{exercise}

% ------------------------------------------------------------------
\section{Several random variables}\label{ch03:sec-multi}

Independence of $X$ and $Y$ means that knowing one gives no information about the other; for densities
it is the factorisation $f_{X,Y}=f_Xf_Y$ \ts{24}{4}{1:03:47}.

\begin{definition}[Covariance, correlation]
$\Cov(X,Y)=\E[(X-\E[X])(Y-\E[Y])]=\E[XY]-\E[X]\,\E[Y]$ (so $\Cov(X,X)=\Var X$) and
\[
  \Corr(X,Y)=\frac{\Cov(X,Y)}{\sqrt{\Var X\,\Var Y}}=\E\Bigl[\frac{X-\mu_X}{\sigma_X}\cdot\frac{Y-\mu_Y}{\sigma_Y}\Bigr]
\]
\ts{24}{4}{1:04:52}\ts{24}{4}{1:05:53}. $X,Y$ are \emph{uncorrelated} if $\Cov(X,Y)=0$.
\end{definition}
Covariance is bilinear; by Cauchy--Schwarz $|\Corr(X,Y)|\le1$, with equality iff $Y=a+bX$ almost surely.
Independent variables are uncorrelated, but not conversely: correlation detects only \emph{linear}
dependence \ts{24}{4}{1:06:58}. Example: $X\sim N(0,1)$ and $Y=X^2$ have $\Cov(X,Y)=\E[X^3]=0$, although $Y$ is a
function of $X$.

\begin{proposition}[Mean and variance of a linear combination]\label{ch03:prop-portvar}
Let $\bm X=(X_1,\dots,X_n)\T$ have mean vector $\bm\mu=\E[\bm X]$ and covariance matrix
$\Sigma=\E[(\bm X-\bm\mu)(\bm X-\bm\mu)\T]$, $\sigma_{ij}=\Cov(X_i,X_j)$ \ts{24}{4}{1:07:59}. For a constant $\bm a\in\R^n$ and
$Y=\bm a\T\bm X$ \ts{24}{4}{1:09:07},
\[
  \E[Y]=\bm a\T\bm\mu,\qquad
  \Var Y=\bm a\T\Sigma\,\bm a=\sum_ia_i^2\Var X_i+2\sum_{i<j}a_ia_j\Cov(X_i,X_j) .
\]
More generally $\E[B\bm X+\bm b]=B\bm\mu+\bm b$ and $\Cov(B\bm X+\bm b)=B\Sigma B\T$.
\end{proposition}
\begin{proof}
$\Var Y=\E[(\bm a\T(\bm X-\bm\mu))^2]=\E[\bm a\T(\bm X-\bm\mu)(\bm X-\bm\mu)\T\bm a]=\bm a\T\Sigma\bm a$: constants come out of the
expectation, which acts entry by entry \ts{24}{4}{1:10:10}. The matrix version is identical.
\end{proof}
In particular $\Sigma$ is positive semi-definite, since $\bm a\T\Sigma\bm a=\Var(\bm a\T\bm X)\ge0$ (\cref{sec:symmetric}).
If $\bm X$ are asset returns and $\bm a$ portfolio weights, $Y$ is the portfolio return: its variance is a
quadratic form in the weights.

\paragraph{Diversification} \ts{24}{4}{1:11:13}\ts{24}{4}{1:12:24}. If the $X_i$ are uncorrelated with common variance
$\sigma^2$ and $a_i=1/n$, then $\sum a_i^2=1/n$ and $\Var Y=\sigma^2/n$ \ts{24}{4}{1:13:27}. With a common correlation
$\rho$ instead (not discussed in class),
\[
  \Var\Bigl(\frac1n\sum_iX_i\Bigr)=\frac{\sigma^2}{n}+\frac{n-1}{n}\,\rho\sigma^2\ \xrightarrow[n\to\infty]{}\ \rho\sigma^2 .
\]

\begin{finance}[Systematic and idiosyncratic risk]
Diversification removes the part of the risk that is specific to each asset, but not the part common to
all of them: with $\rho=0.3$ and $\sigma=30\%$, even an infinitely diversified equal-weighted portfolio keeps a
volatility of $\sqrt{0.3}\cdot30\%\approx16\%$. Only this common, \emph{systematic} risk should command a risk
premium in equilibrium: this is the content of the CAPM (\cref{ch03:sec-capm}) and the starting point of
factor models (\cref{ch:pca}).
\end{finance}

\subsection{Conditional distributions and conditional expectation}\label{ch03:sec-condexp}
This material is only used implicitly in the lectures; it is needed for martingales (\cref{ch:sp1}) and in
2024 PS3 Problem~2.

\begin{definition}[Conditional density and expectation]
If $(X,Y)$ has joint density $f_{X,Y}$, the conditional density of $Y$ given $X=x$ (for $f_X(x)>0$) is
$f_{Y|X}(y\mid x)=f_{X,Y}(x,y)/f_X(x)$, and $m(x)=\E[Y\mid X=x]=\int y\,f_{Y|X}(y\mid x)\dd y$. The
\emph{conditional expectation} $\E[Y\mid X]$ is the random variable $m(X)$. Similarly
$\Var(Y\mid X)=\E[(Y-m(X))^2\mid X]$. (Discrete case: replace densities by pmfs.)
\end{definition}

\begin{proposition}[Properties]\label{ch03:prop-condexp}
For integrable $Y$ and functions $g$:
\begin{enumerate}[(i)]
  \item tower property (law of iterated expectations): $\E\bigl[\E[Y\mid X]\bigr]=\E[Y]$;
  \item taking out what is known: $\E[g(X)Y\mid X]=g(X)\,\E[Y\mid X]$; if $X,Y$ are independent, $\E[Y\mid X]=\E[Y]$;
  \item best prediction: $\E[Y\mid X]$ minimises $\E[(Y-h(X))^2]$ over all functions $h$;
  \item law of total variance: $\Var Y=\E\bigl[\Var(Y\mid X)\bigr]+\Var\bigl(\E[Y\mid X]\bigr)$.
\end{enumerate}
\end{proposition}
\begin{proof}
(i) $\int m(x)f_X(x)\dd x=\iint y\,f_{X,Y}(x,y)\dd y\dd x=\E[Y]$.

(ii) Given $X=x$, $g(X)=g(x)$ is a constant; under
independence $f_{Y|X}=f_Y$.

(iii) $\E[(Y-h)^2]=\E[(Y-m)^2]+\E[(m-h)^2]+2\E[(Y-m)(m-h)(X)]$, and by (i)--(ii) the
cross term is $2\E\bigl[(m-h)(X)\,\E[Y-m(X)\mid X]\bigr]=0$.

(iv) Apply the same decomposition to
$Y-\E[Y]=(Y-m)+(m-\E[Y])$: the cross term vanishes, $\E[(Y-m)^2]=\E[\Var(Y\mid X)]$ by (i), and $\E[m(X)]=\E[Y]$.
\end{proof}

\begin{intuition}[Conditional expectation is a projection]
In the Hilbert space of square-integrable variables with $\langle U,V\rangle=\E[UV]$, $\E[Y\mid X]$ is the orthogonal
projection of $Y$ on the subspace of functions of $X$: (iii) is Pythagoras. Linear regression projects on the
smaller subspace of affine functions $a+bX$ (\cref{ch:regression}); for jointly normal variables the two
projections coincide (\cref{ch03:sec-mvn}). For a general information set $\mathcal G$ (e.g.\ the price history up
to time $t$) $\E[Y\mid\mathcal G]$ is defined by the same orthogonality, $\E[(Y-\E[Y\mid\mathcal G])\,W]=0$ for every
bounded $\mathcal G$-measurable $W$: this is the object in the definition of a martingale.
\end{intuition}

\begin{exercise}[not from the course]\label{ch03:ex-pairwise}
Let $X_1,X_2$ be independent with $\Prob(X_i=1)=\Prob(X_i=-1)=\frac12$, and $X_3=X_1X_2$. Show that $X_1,X_2,X_3$ are pairwise
independent but not mutually independent.
\end{exercise}

% ------------------------------------------------------------------
\section{The multivariate normal distribution}\label{ch03:sec-mvn}

\begin{definition}[Multivariate normal]
$\bm X\in\R^n$ is \emph{multivariate normal}, $\bm X\sim N_n(\bm\mu,\Sigma)$, if every linear combination $\bm a\T\bm X$ is
normal (a constant counting as $N(c,0)$); then $\bm\mu=\E[\bm X]$ and $\Sigma=\Cov\bm X$. Components of such a vector are
called \emph{jointly normal}.
\end{definition}

\begin{proposition}\label{ch03:prop-mvn}
\leavevmode
\begin{enumerate}[(i)]
  \item $\bm X\sim N_n(\bm\mu,\Sigma)$ iff $\E[e^{\bm t\T\bm X}]=\exp(\bm t\T\bm\mu+\frac12\bm t\T\Sigma\bm t)$ for all $\bm t\in\R^n$.
  \item $\bm X\eqd\bm\mu+A\bm Z$ with $\bm Z$ a vector of i.i.d.\ $N(0,1)$ and any $A$ with $AA\T=\Sigma$ (e.g.\ Cholesky, or
        $A=U\Lambda^{1/2}U\T$ from \cref{thm:spectral}).
  \item If $\Sigma$ is invertible, $\bm X$ has density
        $f(\bm x)=(2\pi)^{-n/2}(\det\Sigma)^{-1/2}\exp\bigl(-\frac12(\bm x-\bm\mu)\T\Sigma^{-1}(\bm x-\bm\mu)\bigr)$.
\end{enumerate}
\end{proposition}
\begin{proof}
(i) $\bm t\T\bm X\sim N(\bm t\T\bm\mu,\bm t\T\Sigma\bm t)$; evaluate \eqref{ch03:eq-normalmgf} at $1$. Conversely this MGF restricted to
$\bm t=s\bm a$ is the MGF of $N(\bm a\T\bm\mu,\bm a\T\Sigma\bm a)$.

(ii) $\E[e^{\bm t\T(\bm\mu+A\bm Z)}]=e^{\bm t\T\bm\mu}\prod_i\E[e^{(A\T\bm t)_iZ_i}]
=e^{\bm t\T\bm\mu+\frac12|A\T\bm t|^2}$, and $|A\T\bm t|^2=\bm t\T\Sigma\bm t$; conclude by the $n$-dimensional version of
\cref{ch03:thm-mgf}.

(iii) Change variables $\bm x=\bm\mu+A\bm z$ in the density $(2\pi)^{-n/2}e^{-|\bm z|^2/2}$ of $\bm Z$:
$|\det A|=(\det\Sigma)^{1/2}$ and $|\bm z|^2=(\bm x-\bm\mu)\T\Sigma^{-1}(\bm x-\bm\mu)$.
\end{proof}

\begin{theorem}[Properties of the multivariate normal]\label{ch03:thm-mvn}
Let $\bm X\sim N_n(\bm\mu,\Sigma)$.
\begin{enumerate}[(i)]
  \item $B\bm X+\bm b\sim N(B\bm\mu+\bm b,\,B\Sigma B\T)$; in particular all marginals are normal.
  \item Jointly normal variables are independent iff they are uncorrelated.
  \item Partition $\bm X=(\bm X_1,\bm X_2)$ with $\Sigma_{22}$ invertible. Then
        \[
          \bm X_1\mid\bm X_2=\bm x_2\ \sim\ N\bigl(\bm\mu_1+\Sigma_{12}\Sigma_{22}^{-1}(\bm x_2-\bm\mu_2),\ \Sigma_{11}-\Sigma_{12}\Sigma_{22}^{-1}\Sigma_{21}\bigr).
        \]
\end{enumerate}
\end{theorem}
\begin{proof}
(i) Every linear combination of $B\bm X+\bm b$ is a linear combination of $\bm X$ plus a constant; the moments follow from
\cref{ch03:prop-portvar}.

(ii) If the cross-covariances vanish, the MGF of \cref{ch03:prop-mvn}(i) factorises into the MGFs
of the blocks.

(iii) Let $\bm\varepsilon=\bm X_1-\bm\mu_1-\Sigma_{12}\Sigma_{22}^{-1}(\bm X_2-\bm\mu_2)$. By (i) $(\bm\varepsilon,\bm X_2)$ is
jointly normal, and $\Cov(\bm\varepsilon,\bm X_2)=\Sigma_{12}-\Sigma_{12}\Sigma_{22}^{-1}\Sigma_{22}=0$, so $\bm\varepsilon$ is independent of
$\bm X_2$ by (ii), with mean $\bm0$ and covariance $\Sigma_{11}-2\Sigma_{12}\Sigma_{22}^{-1}\Sigma_{21}+\Sigma_{12}\Sigma_{22}^{-1}\Sigma_{22}\Sigma_{22}^{-1}\Sigma_{21}
=\Sigma_{11}-\Sigma_{12}\Sigma_{22}^{-1}\Sigma_{21}$. Given $\bm X_2=\bm x_2$, $\bm X_1$ is a constant plus $\bm\varepsilon$.
\end{proof}

\noindent\textbf{Bivariate case.} If $(X,Y)$ is jointly normal with means $\mu_X,\mu_Y$, standard deviations
$\sigma,\tau$ and correlation $\rho$, then
\begin{equation}\label{ch03:eq-bivreg}
  Y=a+bX+\varepsilon,\qquad b=\frac{\rho\tau}{\sigma},\quad a=\mu_Y-b\mu_X,\quad \varepsilon\sim N\bigl(0,(1-\rho^2)\tau^2\bigr)\ \text{independent of }X .
\end{equation}
So $\E[Y\mid X]=\mu_Y+\rho\frac{\tau}{\sigma}(X-\mu_X)$ is exactly the regression line, and the conditional variance does
not depend on $X$. This decomposition is the hint of 2024 PS3 Problem~1 (\cref{ch03:ex-stein}).

\paragraph{Geometry.} The level sets of the density are the ellipsoids $(\bm x-\bm\mu)\T\Sigma^{-1}(\bm x-\bm\mu)=c$, centred at
$\bm\mu$, with principal axes along the eigenvectors of $\Sigma$ and semi-axes $\sqrt{c\lambda_i}$; in two dimensions they
are ellipses tilted upward for $\rho>0$ and downward for $\rho<0$ \ts{24}{5}{3:12}\ts{24}{5}{4:20}. Since
$(\bm X-\bm\mu)\T\Sigma^{-1}(\bm X-\bm\mu)=|\bm Z|^2$, this quadratic form has the $\chi^2_n$ law of \cref{ch03:sec-chisq}.

\begin{remark}[Normal marginals are not enough; not discussed in class]
Let $Z\sim N(0,1)$ and $S=\pm1$ with probability $\frac12$ each, independent of $Z$, and $W=SZ$. Then $W\sim N(0,1)$ and
$\Cov(Z,W)=\E[S]\,\E[Z^2]=0$, yet $Z$ and $W$ are dependent ($|W|=|Z|$) and $Z+W=0$ with probability $\frac12$, so
$(Z,W)$ is not jointly normal. Joint normality is a strong assumption about \emph{dependence}, and in
particular it makes simultaneous extreme moves of several assets very unlikely.
\end{remark}

% ------------------------------------------------------------------
\section{Principal components of a random vector}\label{ch03:sec-pca}

\begin{coursenote}[Scope]
This section is the probabilistic part of PCA (slides 26--32 of \file{lec04_1}, \ts{24}{4}{1:13:27} to
\ts{24}{5}{21:04}). Applications to equity indices and to the yield curve \ts{24}{5}{15:51}, announced in L5 for
S.~Andreev's lecture, are in \cref{ch:pca}.
\end{coursenote}

Let $\bm x\in\R^m$ have $\E[\bm x]=\bm\alpha$ and $\Cov\bm x=\Sigma$ \ts{24}{4}{1:14:27}. $\Sigma$ is symmetric and positive
semi-definite \ts{24}{4}{1:15:32}, so by \cref{thm:spectral}
\begin{align*}
  \Sigma&=\Gamma\Lambda\Gamma\T=\sum_{i=1}^m\lambda_i\bm\gamma_i\bm\gamma_i\T,\\
  \Lambda&=\diag(\lambda_1,\dots,\lambda_m),\quad
  \lambda_1\ge\dots\ge\lambda_m\ge0,\quad
  \Gamma=[\bm\gamma_1\cdots\bm\gamma_m],\ \Gamma\T\Gamma=I_m .
\end{align*}

\begin{definition}[Principal component variables]
$\bm p=\Gamma\T(\bm x-\bm\alpha)$, i.e.\ $p_i=\bm\gamma_i\T(\bm x-\bm\alpha)$ \ts{24}{4}{1:17:39}.
\end{definition}

\begin{proposition}\label{ch03:prop-pca}
\leavevmode
\begin{enumerate}[(i)]
  \item $\E[\bm p]=\bm0$ and $\Cov\bm p=\Gamma\T\Sigma\Gamma=\Lambda$: the PC variables are uncorrelated, with variances
        $\lambda_1\ge\dots\ge\lambda_m$ \ts{24}{4}{1:18:39}; and $\bm x=\bm\alpha+\Gamma\bm p$.
  \item Total variance: $\sum_i\Var(x_i)=\tr\Sigma=\tr(\Lambda\Gamma\T\Gamma)=\sum_k\lambda_k=\sum_k\Var(p_k)$ \ts{24}{5}{21:04}.
  \item Variational characterisation \ts{24}{5}{17:53}: $\bm\gamma_1$ maximises $\Var(\bm w\T\bm x)=\bm w\T\Sigma\bm w$ over unit vectors
        $\bm w$; $\bm\gamma_k$ maximises it over unit vectors orthogonal to $\bm\gamma_1,\dots,\bm\gamma_{k-1}$ \ts{24}{5}{18:55}; the
        maxima are $\lambda_1,\dots,\lambda_k$. The last component has the \emph{smallest} variance, and $\lambda_m=0$ iff
        $\bm\gamma_m\T(\bm x-\bm\alpha)=0$ almost surely, an exact linear relation among the $x_i$.
\end{enumerate}
\end{proposition}
\begin{proof}
(i) and (ii) are \cref{ch03:prop-portvar} and $\tr(AB)=\tr(BA)$. (iii) Write $\bm w=\Gamma\bm c$, $|\bm c|=|\bm w|=1$:
$\bm w\T\Sigma\bm w=\sum_i\lambda_ic_i^2\le\lambda_1$, with equality at $\bm c=\bm e_1$. Orthogonality to $\bm\gamma_1,\dots,\bm\gamma_{k-1}$
means $c_1=\dots=c_{k-1}=0$, and then the bound is $\lambda_k$. A variable with zero variance is a.s.\ equal to its mean.
\end{proof}

\begin{intuition}[Shift and rotate]
$\bm x\mapsto\bm p$ moves the origin to the mean and rotates the axes onto the principal axes of the
density ellipsoids; no distance is changed \ts{24}{5}{8:39}. $\Sigma$ is the ``inertia tensor'' of the cloud of
outcomes and PCA diagonalises it. A zero eigenvalue is a flat direction: before the euro the Deutsche Mark and
the French franc moved so closely that the smallest PC of currency returns was almost degenerate
\ts{24}{5}{19:59}.
\end{intuition}

\paragraph{Factor form} \ts{24}{5}{9:43}\ts{24}{5}{10:48}. Split $\Gamma=[\Gamma_1\,\Gamma_2]$ and $\bm p=(\bm p_1,\bm p_2)$ with the $K<m$
largest eigenvalues in the first block. Then
\[
  \bm x=\bm\alpha+\Gamma_1\bm p_1+\Gamma_2\bm p_2=\bm\alpha+B\bm f+\bm\varepsilon,\qquad B=\Gamma_1,\ \bm f=\bm p_1,\ \bm\varepsilon=\Gamma_2\bm p_2 ,
\]
like a factor model with $K$ factors, except that $\Cov\bm\varepsilon=\Gamma_2\Lambda_2\Gamma_2\T$ is not diagonal
\ts{24}{5}{11:48}. If $\sum_{k\le K}\lambda_k/\sum_k\lambda_k$ is close to $1$, the first $K$ components carry almost all the
variability.

\paragraph{Empirical PCA} \ts{24}{5}{12:48}\ts{24}{5}{13:48}. For data $X=[\bm x_1\cdots\bm x_T]\in\R^{m\times T}$ (one column per
date): row means $\bar{\bm x}=\frac1TX\bm1_T$; de-meaned matrix $X^*=X-\bar{\bm x}\bm1_T\T$; sample covariance
$\hat\Sigma=\frac1TX^*X^{*\mathsf T}=\hat\Gamma\hat\Lambda\hat\Gamma\T$; sample PCs $P=\hat\Gamma\T X^*$ ($m\times T$). Equivalently, from the
(reduced) SVD $X^*=VDU\T$ one gets $\hat\Lambda=D^2/T$, $\hat\Gamma=V$, $P=DU\T$ \ts{24}{5}{14:49} (\cref{ch03:ex-pcasvd}).

\begin{coursenote}[Erratum in the slides]
Slide~30 of \file{lec04_1} describes $U$ as an ``$(m\times T)$ row-orthonormal matrix, $UV'=I_m$''. For $X^*=VDU'$ to be
conformable, $U$ must be $T\times m$ with orthonormal columns, $U\T U=I_m$ (equivalently $U\T$ is $m\times T$ with
orthonormal rows); $UV'$ is not even defined unless $T=m$. The formulas $\hat\Lambda=D^2/T$, $\hat\Gamma=V$, $P=DU'$ are
correct (for $T\ge m$).
\end{coursenote}

\begin{exercise}[\file{lec04_1}, slide 30: PCA through the SVD]\label{ch03:ex-pcasvd}
Let $X^*$ be the de-meaned $m\times T$ data matrix and $X^*=VDU\T$ its reduced SVD ($V$ orthogonal $m\times m$, $D$ diagonal
with $d_1\ge\dots\ge d_m\ge0$, $U$ of size $T\times m$ with $U\T U=I_m$). Show that $\hat\Lambda=\frac1TD^2$, $\hat\Gamma=V$ and
$P=\hat\Gamma\T X^*=DU\T$.
\end{exercise}

% ------------------------------------------------------------------
\section{Sampling distributions: \texorpdfstring{$\chi^2$}{chi-square}, \texorpdfstring{$t$}{t} and \texorpdfstring{$F$}{F}}\label{ch03:sec-chisq}

These laws arise from sums of squares of Gaussians and reappear in every test of \cref{ch:regression}
\ts{24}{5}{21:04}.

\begin{definition}[$\chi^2$ distribution]
If $Z\sim N(0,1)$, $U=Z^2$ has the \emph{chi-squared law with 1 degree of freedom}, $\chi^2_1$. If $Z_1,\dots,Z_n$ are
i.i.d.\ $N(0,1)$, $V=Z_1^2+\dots+Z_n^2\sim\chi^2_n$ \ts{24}{5}{22:06}.
\end{definition}
\begin{proposition}
$\chi^2_1=\mathrm{Gamma}(\frac12,\frac12)$ and $\chi^2_n=\mathrm{Gamma}(\frac n2,\frac12)$, with
\[
  f_U(u)=\frac{u^{-1/2}e^{-u/2}}{\sqrt{2\pi}},\qquad
  f_V(v)=\frac{v^{n/2-1}e^{-v/2}}{2^{n/2}\Gamma(n/2)},\qquad M_V(t)=(1-2t)^{-n/2}\ \ (t<\tfrac12),
\]
$\E[V]=n$, $\Var V=2n$.
\end{proposition}
\begin{proof}
$F_U(u)=\Prob(-\sqrt u\le Z\le\sqrt u)=2\Phi(\sqrt u)-1$, so $f_U(u)=\phi(\sqrt u)/\sqrt u$ (the map $z\mapsto z^2$ is two-to-one,
so \cref{ch03:thm-cov} is applied on each branch). For the Gamma$(\alpha,\lambda)$ density,
$\int_0^\infty\frac{\lambda^\alpha}{\Gamma(\alpha)}x^{\alpha-1}e^{-(\lambda-t)x}\dd x=(1-t/\lambda)^{-\alpha}$ for $t<\lambda$, which is
$(1-2t)^{-1/2}$ for $\alpha=\lambda=\frac12$. The MGF of a sum of $n$ independent copies is the $n$-th power, $(1-2t)^{-n/2}$,
the MGF of Gamma$(\frac n2,\frac12)$; conclude by \cref{ch03:thm-mgf}. Mean $\alpha/\lambda=n$, variance $\alpha/\lambda^2=2n$.
\end{proof}

\begin{coursenote}[Erratum in the slides]
Slide~34 of \file{lec04_1} calls $\lambda$ in Gamma$(\alpha,\lambda)$ ``the scale parameter''. With the density
$\lambda^\alpha x^{\alpha-1}e^{-\lambda x}/\Gamma(\alpha)$ used on slide~33, $\lambda$ is the \emph{rate}; the scale is $1/\lambda$ ($=2$ for
$\chi^2$). R uses the same convention: \texttt{dgamma(x, shape = a, rate = l)}.
\end{coursenote}

\begin{definition}[Student $t$ and Fisher $F$]
Let $Z\sim N(0,1)$, $U\sim\chi^2_r$, $V\sim\chi^2_n$, $W\sim\chi^2_m$, independent. Then
$T=Z/\sqrt{U/r}\sim t_r$ (\emph{$t$ distribution with $r$ degrees of freedom}) \ts{24}{5}{23:11} and
$F=\frac{W/m}{V/n}\sim F_{m,n}$ \ts{24}{5}{24:13}.
\end{definition}
Properties:
\begin{itemize}
  \item $t_r$ has density $\frac{\Gamma((r+1)/2)}{\sqrt{r\pi}\,\Gamma(r/2)}(1+x^2/r)^{-(r+1)/2}$: a power-law tail $|x|^{-(r+1)}$, so moments of
        order $\ge r$ are infinite; $\Var T=r/(r-2)$ for $r>2$ and $\kappa=3+6/(r-4)$ for $r>4$. $t_1$ is the Cauchy law,
        and $t_r\to N(0,1)$ as $r\to\infty$ because $U/r\to1$.
  \item $t_r$ has heavier tails than the normal: $Z$ is divided by a random factor that is only \emph{around} $1$
        \ts{24}{5}{24:13}.
  \item $\E[F]=\E[W/m]\,\E[n/V]=n/(n-2)$ for $n>2$ (using $\E[1/V]=1/(n-2)$), and $T^2\sim F_{1,r}$ if $T\sim t_r$.
  \item $F$ statistics compare two independent variance estimates (analysis of variance, regression $F$ tests)
        \ts{24}{5}{25:16}. When the null hypothesis is false, the numerator is a \emph{non-central} $\chi^2$ and the
        statistic has a non-central $F$ law, which governs the power of the test \ts{24}{5}{26:16}.
\end{itemize}
Where they come from (preview of \cref{ch:regression}): for an i.i.d.\ $N(\mu,\sigma^2)$ sample, $(n-1)s^2/\sigma^2\sim\chi^2_{n-1}$
independently of $\bar X$, hence $\sqrt n(\bar X-\mu)/s\sim t_{n-1}$ (Student, 1908).

\subsection{Heavy tails}\label{ch03:sec-heavy}

\begin{figure}[ht]
\centering
\begin{tikzpicture}
\begin{axis}[width=0.8\linewidth,height=5.6cm,ymode=log,xmin=0,xmax=6,ymin=1e-9,ymax=1,
  axis lines=left,xlabel={$x$ (in standard deviations)},ylabel={density (log scale)},
  legend style={font=\scriptsize,draw=none,fill=none},legend pos=south west,
  tick label style={font=\small},label style={font=\small}]
  \addplot[very thick,defcol,domain=0:6,samples=121] {exp(-x^2/2)/sqrt(2*pi)};
  \addlegendentry{$N(0,1)$}
  \addplot[very thick,thmcol,dashed,domain=0:6,samples=121] {0.49007*(1+x^2/3)^(-3)};
  \addlegendentry{$t_5$ rescaled to variance 1}
\end{axis}
\end{tikzpicture}
\caption{Right tails of the standard normal and of a Student $t_5$ scaled to unit variance (density
$0.490\,(1+x^2/3)^{-3}$, kurtosis 9). On a log scale the Gaussian tail is a parabola, the $t$ tail a
logarithm.}\label{ch03:fig-tails}
\end{figure}

With equal variances the two laws of \cref{ch03:fig-tails} look alike near the centre, but
$\Prob(|X|>4)$ is $6.3\times10^{-5}$ for the normal and $3.6\times10^{-3}$ for the rescaled $t_5$, about 56 times more.
For daily returns with 252 trading days a year, a 4-standard-deviation day happens on average once every
63~years under the normal model and about once a year under the $t_5$ model.

\begin{proposition}[Normal variance mixtures; not discussed in class]\label{ch03:prop-mixture}
Let $X\mid V\sim N(\mu,\sigma^2V)$ with $V>0$ random and $\E[V^2]<\infty$. Then $\E[X]=\mu$, $\Var X=\sigma^2\E[V]$ and
\[
  \kappa_X=3\,\frac{\E[V^2]}{(\E[V])^2}=3\Bigl(1+\frac{\Var V}{(\E[V])^2}\Bigr)\ \ge\ 3,
\]
with equality iff $V$ is constant.
\end{proposition}
\begin{proof}
Condition on $V$: $\E[(X-\mu)^2\mid V]=\sigma^2V$ and $\E[(X-\mu)^4\mid V]=3\sigma^4V^2$; take expectations (tower property).
\end{proof}
The $t_r$ law is such a mixture: $T=Z\sqrt{r/U}$, so $T\mid U\sim N(0,r/U)$. \emph{Stochastic volatility} models, in
which the variance of returns is itself random, are of this type, and always produce kurtosis above 3. 2024
PS3 Problem~2 extends Stein's lemma (\cref{ch03:sec-capm}) to them (\cref{ch03:ex-sv}).

\begin{intuition}[Superstatistics]
A Gaussian whose variance fluctuates is like a Brownian particle in a medium whose temperature fluctuates
slowly: locally Gaussian, globally fat-tailed. Physicists call this superposition of statistics
\emph{superstatistics}; in finance the fluctuating temperature is the volatility, which clusters in time
(\cref{ch:volatility,ch:timeseries}).
\end{intuition}

\begin{finance}[Heavy tails in practice]
Daily returns of stocks and indices have sample kurtosis far above 3, and many studies find tails decaying
roughly like a power law, with exponent near 3 for individual stocks (the ``inverse cubic law''). Consequences:
normal-based value-at-risk understates extreme losses (\cref{ch:var}); sample moments of high order are
unstable or meaningless; averages converge more slowly than Gaussian intuition suggests, and not at all for a
Cauchy-like law (\cref{ch03:ex-cauchy}).
\end{finance}

\begin{exercise}[not from the course: a heavy-tailed mixture]
Let $X\mid V\sim N(0,V)$ with $V$ exponential with rate 1. Show that $M_X(t)=1/(1-t^2/2)$ for $|t|<\sqrt2$, identify the law of
$X$ (Laplace, density $\frac1{\sqrt2}e^{-\sqrt2|x|}$), and check \cref{ch03:prop-mixture}: $\Var X=1$, $\kappa=6$.
\end{exercise}

% ------------------------------------------------------------------
\section{Limit theorems}\label{ch03:sec-limits}

\begin{definition}[Modes of convergence]
$X_n\to X$ \emph{in probability} if $\Prob(|X_n-X|>\varepsilon)\to0$ for every $\varepsilon>0$; \emph{in distribution} if
$F_{X_n}(x)\to F_X(x)$ at every continuity point of $F_X$; \emph{almost surely} if $\Prob(X_n\to X)=1$. Almost sure
convergence implies convergence in probability, which implies convergence in distribution.
\end{definition}

\begin{lemma}[Markov and Chebyshev inequalities]\label{ch03:lem-cheb}
For $a>0$: $\Prob(|X|\ge a)\le\E[|X|]/a$. If $\Var X=\sigma^2<\infty$ and $\varepsilon>0$:
$\Prob(|X-\mu|\ge\varepsilon)\le\sigma^2/\varepsilon^2$.
\end{lemma}
\begin{proof}
$a\,\ind\{|X|\ge a\}\le|X|$ pointwise; take expectations. For Chebyshev apply Markov to $(X-\mu)^2$ and $a=\varepsilon^2$: the
parabola $(x-\mu)^2$ lies above the step function $\varepsilon^2\ind\{|x-\mu|\ge\varepsilon\}$ (\cref{ch03:fig-cheb}), and averaging
preserves the inequality \ts{24}{5}{29:29}\ts{24}{5}{30:30}.
\end{proof}

\begin{figure}[ht]
\centering
\begin{tikzpicture}
\begin{axis}[width=0.62\linewidth,height=4.8cm,axis lines=middle,xmin=-2.3,xmax=2.3,ymin=0,ymax=4.4,
  xtick={-1,1},xticklabels={$-\varepsilon$,$\varepsilon$},ytick={1},yticklabels={$\varepsilon^2$},
  xlabel={$d$},x label style={anchor=west},tick label style={font=\small}]
  \addplot[very thick,defcol,domain=-2.05:2.05,samples=83] {x^2};
  \draw[very thick,thmcol] (axis cs:-2.2,1) -- (axis cs:-1,1);
  \draw[very thick,thmcol] (axis cs:-1,0.02) -- (axis cs:1,0.02);
  \draw[very thick,thmcol] (axis cs:1,1) -- (axis cs:2.2,1);
  \draw[dashed,black!40] (axis cs:-1,0) -- (axis cs:-1,1);
  \draw[dashed,black!40] (axis cs:1,0) -- (axis cs:1,1);
  \node[defcol,font=\small,anchor=west] at (axis cs:1.55,3.6) {$d^2$};
  \node[thmcol,font=\small,anchor=south] at (axis cs:-1.6,1.05) {$\varepsilon^2\ind\{|d|\ge\varepsilon\}$};
\end{axis}
\end{tikzpicture}
\caption{The picture behind Chebyshev's inequality, as drawn in class \ts{24}{5}{29:29}: with $d=\bar X_n-\mu$,
$d^2\ge\varepsilon^2\ind\{|d|\ge\varepsilon\}$, so $\Var\bar X_n=\E[d^2]\ge\varepsilon^2\Prob(|d|\ge\varepsilon)$.}\label{ch03:fig-cheb}
\end{figure}

\begin{theorem}[Weak law of large numbers]\label{ch03:thm-wlln}
Let $X_1,X_2,\dots$ be i.i.d.\ with $\E[X_i]=\mu$ and $\Var X_i=\sigma^2<\infty$, and $\bar X_n=\frac1n(X_1+\dots+X_n)$. Then
$\bar X_n\to\mu$ in probability \ts{24}{5}{27:20}\ts{13}{3}{49:04}:
$\Prob(|\bar X_n-\mu|>\varepsilon)\le\dfrac{\sigma^2}{n\varepsilon^2}\to0$.
\end{theorem}
\begin{proof}
$\E[\bar X_n]=\mu$ by linearity. $\Var\bar X_n=\frac1{n^2}\sum_{i,j}\Cov(X_i,X_j)=\frac{\sigma^2}n$, since the cross terms vanish
\ts{13}{3}{58:43}. Apply \cref{ch03:lem-cheb} \ts{24}{5}{28:28}.
\end{proof}

\begin{remark}
\leavevmode
\begin{itemize}
  \item The proof only uses that the $X_i$ are pairwise \emph{uncorrelated} with equal means and variances.
  \item The conclusion holds with only $\E[|X_i|]<\infty$ (Khinchin's WLLN) \ts{24}{5}{31:38}, and the \emph{strong}
        law gives $\bar X_n\to\mu$ almost surely under the same assumption \ts{13}{3}{1:02:53}.
  \item Without a finite mean it fails: the mean of $n$ i.i.d.\ Cauchy variables is again Cauchy (\cref{ch03:ex-cauchy}).
\end{itemize}
\end{remark}

\begin{example}[How many observations?]\label{ch03:ex-samplesize}
To be 99\% sure that $|\bar X_n-\mu|<\varepsilon=0.1$
\ts{13}{3}{1:00:50}, Chebyshev requires
$\sigma^2/(n\varepsilon^2)\le0.01$; with $\sigma=1$ (our choice) this
is $n\ge10{,}000$. The CLT below suggests
$n\approx(2.576\,\sigma/\varepsilon)^2\approx664$.\sloppy Chebyshev is valid
for every distribution with finite variance, hence crude; sharper
exponential bounds (Chernoff, Hoeffding) use MGF-type arguments, as
mentioned in class \ts{13}{3}{1:01:53}.
\end{example}

\begin{finance}[The casino, the poker room and the trader (2013)]
Casino games against the house are designed so that the house has an edge of about 1--5\% (2013 notes)
while the variance of each game is large compared with the edge \ts{13}{3}{50:22}. If a player wins an
even-money bet with probability $0.48$ (the indicative figure quoted in class for blackjack
\ts{13}{3}{51:22}), the expected P\&L per \$1 bet is $0.48-0.52=-\$0.04$ with a standard deviation of about
\$1: over an evening the edge is invisible, but the casino plays an enormous $n$ and the LLN makes its
profit almost certain \ts{13}{3}{52:22}. In poker the players play against each other and the house simply
charges a fee per hand (the \emph{rake}) \ts{13}{3}{54:32}; a player whose edge over the others exceeds the rake
can let the LLN work for him. The same logic drives a hedge fund or a high-frequency trader: a small edge
per trade, many (nearly) independent trades \ts{13}{3}{55:32}. The practical difficulty is the variance:
drawdowns make it hard to tell a real edge from bad luck \ts{13}{3}{56:34}, and correlated trades reduce the
effective $n$.
\end{finance}

\begin{theorem}[Central limit theorem]\label{ch03:thm-clt}
Let $X_1,X_2,\dots$ be i.i.d.\ with mean $\mu$ and variance $\sigma^2\in(0,\infty)$. Then
\[
  Z_n=\frac{1}{\sqrt n}\sum_{i=1}^n(X_i-\mu)=\sqrt n\,(\bar X_n-\mu)\ \longrightarrow\ N(0,\sigma^2)\quad\text{in distribution}
\]
\ts{24}{5}{31:38}\ts{13}{3}{1:07:29}.
\end{theorem}
\begin{proof}[Proof when the MGF exists]
Both courses prove it this way \ts{24}{5}{33:48}\ts{13}{3}{1:09:46}. Take $\mu=0$ and let $M$ be the MGF of $X_i$, finite near $0$.
By independence the expectation of the product is the product of expectations \ts{24}{5}{34:51}:
\[
  M_{Z_n}(t)=\E\Bigl[\prod_{i=1}^ne^{tX_i/\sqrt n}\Bigr]=M\Bigl(\frac{t}{\sqrt n}\Bigr)^n .
\]
$M$ is smooth near $0$ with $M(0)=1$, $M'(0)=\E[X]=0$, $M''(0)=\E[X^2]=\sigma^2$, so by Taylor
$M(s)=1+\frac12\sigma^2s^2+o(s^2)$ as $s\to0$ \ts{24}{5}{35:54}. For fixed $t$,
\[
  \ln M_{Z_n}(t)=n\ln\Bigl(1+\frac{\sigma^2t^2}{2n}+o\bigl(\tfrac1n\bigr)\Bigr)\ \longrightarrow\ \frac{\sigma^2t^2}{2},
\]
as in $(1+r/n)^n\to e^r$, the limit behind continuous compounding \ts{24}{5}{37:01}. So $M_{Z_n}(t)\to e^{\sigma^2t^2/2}$, the MGF
of $N(0,\sigma^2)$, and \cref{ch03:thm-mgf}(ii) concludes.
\end{proof}

\begin{remark}
\leavevmode
\begin{itemize}
  \item The statement ``$\bar X_n\approx N(\mu,\sigma^2/n)$'' is written with quotes in class, because a limit law cannot
        depend on $n$; it is the practical approximation \ts{24}{5}{32:47}.
  \item The MGF is only a convenience: with characteristic functions the same proof works whenever $\sigma^2<\infty$.
  \item Identical distribution is not needed (Lindeberg--Lyapunov CLT) \ts{24}{5}{38:04}, and weak dependence is
        allowed (mixing sequences, martingale differences).
\end{itemize}
\end{remark}

To a student's question whether the CLT is of any use when returns
are neither independent nor identically distributed \ts{24}{5}{38:04},
the answer in class: what matters is that a quantity be a sum of many
small, roughly independent contributions. Projections of a
high-dimensional vector on a random direction, and the residuals of a
regression, are such sums and are often close to Gaussian
\ts{24}{5}{39:09}\ts{24}{5}{40:17}.

\begin{example}[Estimating a mean, 2013 notes Example 4.2]
To estimate an unknown mean from i.i.d.\ observations one uses
$\bar X_n$ \ts{13}{3}{1:16:15}: by the LLN it converges to $\mu$, and by
the CLT its error is approximately $N(0,\sigma^2/n)$ \ts{13}{3}{1:17:21},
which gives the usual interval $\bar X_n\pm1.96\,s/\sqrt n$.
\end{example}

\begin{coursenote}[Is the sample mean the best estimator?]
In the 2013 lecture the sample mean is called the maximum likelihood estimator \ts{13}{3}{1:18:23}. It is the MLE
for normal (and Poisson, exponential) data but not in general, and the notes rightly add that for some
distributions it is not the best choice: for Cauchy data it does not converge at all, while the sample
median does; for heavy-tailed data robust estimators are preferable (\cref{ch:regression}).
\end{coursenote}

\begin{finance}[Aggregation of returns]
A $T$-day log-return is a sum of $T$ daily log-returns, so by the CLT it is approximately normal for long
horizons even when daily returns are fat-tailed (as long as their variance is finite): return distributions
look more Gaussian at monthly than at daily frequency. Volatility clustering slows this convergence down. The
same theorem makes the binomial tree converge to the lognormal model (\cref{ch:bs}) and the random walk to
Brownian motion (\cref{ch:sp1}).
\end{finance}

\begin{exercise}[not from the course: no LLN for Cauchy]\label{ch03:ex-cauchy}
Assume that the standard Cauchy law has characteristic function $e^{-|t|}$. Show that if $X_1,\dots,X_n$ are i.i.d.\
standard Cauchy, $\bar X_n$ is again standard Cauchy. Why does this not contradict \cref{ch03:thm-wlln}? What does it
suggest about averaging returns with very heavy tails?
\end{exercise}

% ------------------------------------------------------------------
\section{Probability theory for asset pricing: from expected utility to the CAPM}\label{ch03:sec-capm}

The slides \file{lec04_2} derive the capital asset pricing model with nothing more than this chapter:
normal payoffs, covariances and a Gaussian identity \ts{24}{5}{40:17}. The development follows F.~de Jong and
B.~Rindi, \emph{The Microstructure of Financial Markets} (2009) \ts{24}{5}{1:11:22}.

\subsection{One risky asset}
One period, from $t_0$ (trading) to $t_1$ (end of period) \ts{24}{5}{41:21}. Before $t_0$ an agent is endowed
with $Q_a$ shares of a risky asset and $Q_f$ units of a risk-free asset. At $t_0$ the risk-free asset costs \$1 per
unit and the risky asset $p_a$ per share \ts{24}{5}{42:26}, so the initial wealth is $w_0=Q_ap_a+Q_f$. At $t_1$ the
risk-free asset pays $1+r_f$ per unit and the risky one pays $\tilde F\sim N(\mu,\sigma^2)$ per share \ts{24}{5}{43:28}.
(A normal payoff can be negative; the model is useful if this has negligible probability \ts{24}{5}{44:28}.)
The agent buys $X$ shares at $t_0$ ($X<0$: sells), financed with the risk-free asset, so the final wealth is
\ts{24}{5}{45:31}
\begin{equation}\label{ch03:eq-wealth}
  \tilde w=(Q_a+X)\tilde F+(Q_f-Xp_a)(1+r_f),\qquad \frac{\partial\tilde w}{\partial X}=\tilde F-p_a(1+r_f).
\end{equation}

\paragraph{Expected utility.} The agent chooses $X$ to maximise $\E[U(\tilde w)]$ for a utility function with
$U'>0$ \ts{24}{5}{46:37} (more wealth is better) and $U''<0$ \ts{24}{5}{47:37} (decreasing marginal utility). Why not
simply $U(w)=w$ \ts{24}{5}{48:39}? Because then there is no penalty for risk \ts{24}{5}{49:47}: with $U$ concave,
Jensen's inequality gives $\E[U(\tilde w)]\le U(\E[\tilde w])$, so the agent prefers a sure amount to a gamble with the
same mean. The ratio $A(w)=-U''(w)/U'(w)$ is the (Arrow--Pratt) coefficient of absolute risk aversion.

\begin{coursenote}[Shannon, Kelly and the log-utility debate]
The lecturer recalls a controversy between Claude Shannon and (he thinks) Milton Friedman on maximising the
expected logarithm of wealth \ts{24}{5}{49:47}. The well-documented debate is between Paul Samuelson, who argued
against maximising expected log-wealth, and the proponents of the \emph{Kelly criterion} (J.~Kelly,
C.~Shannon, E.~Thorp); see W.~Poundstone, \emph{Fortune's Formula} (2005). Log utility has $A(w)=1/w$.
\end{coursenote}

\subsection{Stein's lemma}
\begin{lemma}[Stein]\label{ch03:lem-stein}
Let $Y$ and $Y^*$ be jointly normal with $Y\sim N(\mu,\sigma^2)$, and let $g$ be differentiable with
$\E[|g(Y)(Y-\mu)|]<\infty$ and $\E[|g'(Y)|]<\infty$. Then \ts{24}{5}{50:58}\ts{24}{5}{52:01}
\[
  \E[g(Y)(Y-\mu)]=\sigma^2\,\E[g'(Y)]\qquad\text{and}\qquad \Cov\bigl(g(Y),Y^*\bigr)=\E[g'(Y)]\,\Cov(Y,Y^*).
\]
\end{lemma}
\noindent The proof is 2024 PS3 Problem~1 (\cref{ch03:ex-stein}). The first identity is a Gaussian integration
by parts, the second reduces to the first through the decomposition \eqref{ch03:eq-bivreg}.

\begin{intuition}[Gaussian integration by parts]
For a centred Gaussian, $\E[Y\,g(Y)]=\sigma^2\E[g'(Y)]$ is the identity behind Wick's theorem and the
Schwinger--Dyson equations: with $g(y)=y^{k-1}$ it gives $\E[Y^k]=(k-1)\sigma^2\E[Y^{k-2}]$, hence $\E[Y^4]=3\sigma^4$. The second
form says that the covariance of a nonlinear function $g(Y)$ with any jointly Gaussian $Y^*$ is that of its
linearisation, $g'\cdot\Cov(Y,Y^*)$, with the slope $g'$ averaged over the law of $Y$. In the pricing argument
$g=U'$: it converts ``covariance of marginal utility with the payoff'' into ``covariance of wealth with the
payoff''.
\end{intuition}

\begin{exercise}[PS3 (2024), Problem 1: Stein's lemma]\label{ch03:ex-stein}
Prove \cref{ch03:lem-stein} (with $X$ in place of $Y$ and $Y$ in place of $Y^*$): if $X,Y$ are jointly normal with
$X\sim N(\mu,\sigma^2)$ and $g$ is differentiable with $\E[g(X)(X-\mu)]$ and $\E[g'(X)]$ finite, then
$\E[g(X)(X-\mu)]=\sigma^2\E[g'(X)]$ and $\Cov(g(X),Y)=\E[g'(X)]\Cov(X,Y)$.
\emph{Hint:} integrate by parts; you may use that for jointly normal $(X,Y)$ with standard deviations $\sigma,\tau$ and
correlation $\rho$, $Y=a+bX+\varepsilon$ with $a=\mu_Y-b\mu_X$, $b=\rho\tau/\sigma$ and $\varepsilon\sim N(0,(1-\rho^2)\tau^2)$ independent of $X$.
\end{exercise}

\begin{exercise}[PS3 (2024), Problem 2: Stein's lemma for stochastic volatility]\label{ch03:ex-sv}
(From A.~Gron, B.~J{\o}rgensen, N.~Polson, \emph{Optimal portfolio choice and stochastic volatility}, Applied
Stochastic Models in Business and Industry, 2012; cited as a 2011 working paper in the problem set.)
$X$ has \emph{stochastic volatility} if $X\mid V\sim N(\mu,\sigma^2V)$ with $V>0$ random with density $p(V)$; the
law of $X$ is $p(X)=\int_0^\infty p(X\mid V)p(V)\dd V$ and has heavy tails. Assume $0<\E[V]<\infty$ and define the
\emph{size-biased} density $q(V)=V\,p(V)/\E[V]$ and the induced $q(X)=\int_0^\infty p(X\mid V)q(V)\dd V$. The
authors' theorem states
\[
  \Cov\bigl(g(X),X\bigr)=\E^Q[g'(X)]\,\Var X,\qquad(**)
\]
with $\E^Q$ the expectation under the measure induced by $q(V)$. Their proof reads
\begin{align*}
  \Cov(g(X),X)&\overset{(1)}{=}\E\bigl[g(X)(X-\E[X])\bigr]\\
  &\overset{(2)}{=}\E_V\bigl[\E_{X|V}[g(X)(X-\E[X])]\bigr]\\
  &\overset{(3)}{=}\E_V\bigl[\E_{X|V}[g(X)(X-\E[X\mid V])]\bigr].
\end{align*}
\begin{enumerate}[(a)]
  \item Show that $q$ is a probability density on $(0,\infty)$.
  \item Prove (1).
  \item Does (2) follow from the law of
      iterated expectations?
  \item Prove (3) (compare $\E[X\mid V]$ and $\E[X]$).
  \item Applying Stein's lemma conditionally
      on $V$ they obtain $\Cov(g(X),X)=\E_V\{\E_{X|V}[g'(X)]\,\sigma^2V\}=\E\bigl[g'(X)\tfrac{V}{\E[V]}\bigr]\sigma^2\E[V]$ $(*)$. Does $(*)$
      imply $(**)$? Explain.
\end{enumerate}
\end{exercise}

\subsection{First-order condition and the equilibrium price}
Assume one may differentiate under the expectation, which holds by
dominated convergence if, e.g., $U'$ is bounded, or for the exponential
utility below since normal variables have exponential moments
\ts{24}{5}{53:03}\ts{24}{5}{54:16}. By \eqref{ch03:eq-wealth} the first-order
condition is
\begin{align*}
  0&=\frac{\partial}{\partial X}\E[U(\tilde w)]
  =\E\bigl[U'(\tilde w)\bigl(\tilde F-p_a(1+r_f)\bigr)\bigr] \\
   &=\Cov\bigl(U'(\tilde w),\tilde F\bigr)
   +\E[U'(\tilde w)]\bigl(\mu-p_a(1+r_f)\bigr),
\end{align*}
using $\E[AB]=\Cov(A,B)+\E[A]\,\E[B]$ and that a constant shift does not
change a covariance \ts{24}{5}{55:19}. Now $\tilde w$ is normal (affine
in $\tilde F$), so \cref{ch03:lem-stein} with $g=U'$, $Y=\tilde w$,
$Y^*=\tilde F$ gives
\begin{align*}
  \Cov(U'(\tilde w),\tilde F)&=\E[U''(\tilde w)]\Cov(\tilde w,\tilde F)\\
  &=\E[U''(\tilde w)](Q_a+X)\sigma^2
\end{align*}
\ts{24}{5}{56:22}\ts{24}{5}{57:28}. Solving for $p_a$:
\begin{equation}\label{ch03:eq-price1}
  p_a=\frac{1}{1+r_f}\Bigl[\E[\tilde F]
  +\frac{\E[U''(\tilde w)]}{\E[U'(\tilde w)]}\,(Q_a+X)\Var(\tilde F)\Bigr].
\end{equation}
Since $U''<0<U'$ the second term is negative: the price is the
discounted expected payoff \emph{minus a risk adjustment}
\ts{24}{5}{58:29}.

\paragraph{CARA utility} \ts{24}{5}{59:30}. For $U(w)=c-e^{-Aw}$ with
$A>0$: $U'(w)=Ae^{-Aw}>0$ and $U''(w)=-A^2e^{-Aw}=-A\,U'(w)$, so the
ratio in \eqref{ch03:eq-price1} is the constant $-A$ and
\begin{equation}\label{ch03:eq-cara}
  p_a=\frac{1}{1+r_f}\Bigl[\E[\tilde F]-A\,(Q_a+X)\Var(\tilde F)\Bigr].
\end{equation}
The risk adjustment is proportional to the risk aversion $A$, to the
position held $Q_a+X$ and to the variance of the payoff
\ts{24}{5}{1:00:36}.

\begin{coursenote}[Erratum in the slides]
Slide~6 of \file{lec04_2} writes $U'(\tilde w)=AU(\tilde w)$ and $U''(\tilde w)=-A^2U(\tilde w)$. For $U=c-e^{-Aw}$ one has
$U'=Ae^{-Aw}=A(c-U)$ and $U''=-A^2e^{-Aw}$; even with $c=0$ the signs are wrong ($U'=-AU$, $U''=A^2U$). What the
derivation uses, $U''/U'=-A$, and hence \eqref{ch03:eq-cara}, are correct.
\end{coursenote}

\begin{remark}[CARA and normal = mean--variance; not discussed in class]\label{ch03:rem-meanvar}
Since $-A\tilde w$ is normal, \eqref{ch03:eq-normalmgf} at $t=1$ gives $\E[U(\tilde w)]=c-\exp\bigl(-A\,\E[\tilde w]+\frac{A^2}{2}\Var\tilde w\bigr)$.
Maximising expected utility is therefore equivalent to maximising the \emph{certainty equivalent}
$\E[\tilde w]-\frac A2\Var\tilde w$, a mean--variance criterion (\cref{ch:portfolio}). Its maximiser is the demand
$Q_a+X=\bigl(\E[\tilde F]-(1+r_f)p_a\bigr)/(A\sigma^2)$, and \eqref{ch03:eq-cara} is this demand curve solved for the price:
the price at which the agent is content to hold $Q_a+X$ shares.
\end{remark}

\paragraph{Expected return.} With $\tilde r=(\tilde F-p_a)/p_a$, \eqref{ch03:eq-price1} gives
\[
  \E[\tilde r]=r_f-\frac{\E[U''(\tilde w)]}{\E[U'(\tilde w)]}\,\frac{(Q_a+X)\Var(\tilde F)}{p_a}
  \ \overset{\text{CARA}}{=}\ r_f+A\,\frac{(Q_a+X)\Var(\tilde F)}{p_a}:
\]
the risk-free rate plus a risk premium that grows with risk aversion, position and variance.

\subsection{Many risky assets}
Now the agent holds $\bm Q\in\R^n$ shares of $n$ risky assets priced $\bm p$ at $t_0$, whose payoffs
$\tilde{\bm F}\sim N_n(\bm\mu_F,\Sigma_F)$ are jointly normal \ts{24}{5}{1:01:39}\ts{24}{5}{1:02:45}, and buys $\bm X$:
$\tilde w=(\bm Q+\bm X)\T\tilde{\bm F}+(Q_f-\bm X\T\bm p)(1+r_f)$, again normal. The $n$ first-order conditions are handled
exactly as before, with Stein's lemma applied to the jointly normal pair $(\tilde w,\tilde F_i)$:
\[
  p_i=\frac{1}{1+r_f}\Bigl[\E[\tilde F_i]+\frac{\E[U''(\tilde w)]}{\E[U'(\tilde w)]}\Cov(\tilde w,\tilde F_i)\Bigr],\qquad
  \Cov(\tilde w,\tilde F_i)=\bigl[\Sigma_F(\bm Q+\bm X)\bigr]_i .
\]
For CARA, $\bm p=\bigl[\bm\mu_F-A\,\Sigma_F(\bm Q+\bm X)\bigr]/(1+r_f)$. With $\tilde r_i=(\tilde F_i-p_i)/p_i$, so that
$\Cov(\tilde w,\tilde r_i)=\Cov(\tilde w,\tilde F_i)/p_i$,
\[
  \E[\tilde r_i]=r_f-\frac{\E[U''(\tilde w)]}{\E[U'(\tilde w)]}\Cov(\tilde w,\tilde r_i)\ \overset{\text{CARA}}{=}\ r_f+A\Cov(\tilde w,\tilde r_i)
\]
\ts{24}{5}{1:03:48}. The premium of asset $i$ is set by its covariance with the agent's \emph{total} wealth, not by its
own variance: an asset that hedges wealth ($\Cov<0$) is expected to earn less than $r_f$.

\subsection{The capital asset pricing model}
\paragraph{Market portfolio.} In equilibrium prices must be such that the assets are willingly held. With
identical agents (or a representative agent) holding no initial risky endowment, $\bm Q=\bm0$, the agent's
equilibrium purchase $\bm X$ is the whole supply of risky assets: the \emph{market portfolio}
\ts{24}{5}{1:03:48}\ts{24}{5}{1:04:53}. Its payoff, price and return are
\[
  \tilde F_m=\sum_iX_i\tilde F_i,\qquad p_m=\sum_iX_ip_i,\qquad
  \tilde r_m=\frac{\tilde F_m-p_m}{p_m}=\sum_i\omega_i\tilde r_i,\quad \omega_i=\frac{X_ip_i}{p_m},
\]
with $\omega_i$ the market-value weights \ts{24}{5}{1:05:56}. The key identity is
$\tilde w=\tilde F_m+(Q_f-p_m)(1+r_f)=p_m\tilde r_m+\text{const}$, hence $\Cov(\tilde w,\cdot)=p_m\Cov(\tilde r_m,\cdot)$.

\begin{theorem}[CAPM, security market line]\label{ch03:thm-capm}
In this equilibrium, with $\beta_i=\Cov(\tilde r_i,\tilde r_m)/\Var(\tilde r_m)$,
\[
  \E[\tilde r_i]=r_f+\beta_i\bigl(\E[\tilde r_m]-r_f\bigr),\qquad
  p_i=\frac{\E[\tilde F_i]}{1+r_f+\beta_i\bigl(\E[\tilde r_m]-r_f\bigr)} .
\]
\end{theorem}
\begin{proof}
Let $\theta=-\E[U''(\tilde w)]/\E[U'(\tilde w)]>0$. Averaging the expected returns with weights $\omega_i$
\ts{24}{5}{1:07:04}: $\E[\tilde r_m]=r_f+\theta\Cov(\tilde w,\tilde r_m)=r_f+\theta\,p_m\Var(\tilde r_m)$, so
$\theta=(\E[\tilde r_m]-r_f)/(p_m\Var\tilde r_m)$. Then for each asset
$\E[\tilde r_i]-r_f=\theta\Cov(\tilde w,\tilde r_i)=\theta\,p_m\Cov(\tilde r_m,\tilde r_i)=\beta_i(\E[\tilde r_m]-r_f)$ \ts{24}{5}{1:08:10}. Finally
$\E[\tilde r_i]=\E[\tilde F_i]/p_i-1$ gives the price \ts{24}{5}{1:10:18}.
\end{proof}

\begin{figure}[ht]
\centering
\begin{tikzpicture}
\begin{axis}[width=0.62\linewidth,height=5cm,axis lines=left,xmin=0,xmax=2.1,ymin=0,ymax=15,
  xlabel={$\beta_i$},ylabel={$\E[\tilde r_i]$ (\%)},xtick={0,0.5,1,1.5,2},ytick={2,8,14},
  tick label style={font=\small},label style={font=\small}]
  \addplot[very thick,defcol,domain=0:2,samples=2] {2+6*x};
  \addplot[only marks,mark=*,mark size=2pt,thmcol] coordinates {(0,2) (1,8)};
  \node[font=\scriptsize,anchor=west] at (axis cs:0.05,1.3) {$r_f$};
  \node[font=\scriptsize,anchor=north west] at (axis cs:1.02,8) {market portfolio};
  \node[font=\scriptsize,anchor=south east,defcol] at (axis cs:1.9,13.6) {SML};
  \addplot[only marks,mark=o,mark size=2pt,intcol] coordinates {(0.6,8)};
  \node[font=\scriptsize,anchor=south,intcol] at (axis cs:0.6,8.3) {$\alpha>0$};
\end{axis}
\end{tikzpicture}
\caption{Security market line with illustrative numbers $r_f=2\%$, $\E[\tilde r_m]=8\%$. In the CAPM every asset lies
on the line; an asset above it (open dot) would have a positive ``alpha''.}\label{ch03:fig-sml}
\end{figure}

\begin{coursenote}[Notation and assumptions in \file{lec04_2}]
The slides denote the market weights by $\mu_i$, which clashes with the means; here they are $\omega_i$. The step
``the agent's portfolio is the market portfolio'' is the market-clearing condition with identical agents,
stated on slide~12 as an assumption. The CAPM was formulated by W.~Sharpe (1964) and J.~Lintner (1965) from
mean--variance portfolio choice \ts{24}{5}{1:11:22}; by \cref{ch03:rem-meanvar} the CARA-normal route taken here is
equivalent. A Lintner paper was distributed on Canvas; the OCW Week~3 page lists both Lintner (1965) papers.
\end{coursenote}

\begin{finance}[Reading and using the CAPM]
\begin{itemize}
  \item $\beta_i$ is the slope of the regression of $\tilde r_i$ on $\tilde r_m$ \ts{24}{5}{1:09:14} (\cref{def:alphabeta},
        \cref{ch:regression}). High-beta stocks must offer higher expected returns; $\beta<1$ is ``defensive''.
  \item Only systematic risk is priced: the idiosyncratic part of $\Var\tilde r_i$ can be diversified away
        (\cref{ch03:sec-multi}) and earns nothing.
  \item Practice (not discussed in class): the cost of equity in corporate valuation is often set to
        $r_f+\beta\times$(equity risk premium); fund performance is measured by the \emph{alpha} over the SML.
        Empirically the SML is flatter than predicted (low-beta stocks earn more than the CAPM says), which
        motivated multi-factor models (\cref{ch:pca,ch:portfolio}).
\end{itemize}
\end{finance}

% ------------------------------------------------------------------
\begin{keyformulas}
\begin{align*}
&\text{PIT:} && F_X(X)\sim U(0,1)\ (F_X\text{ continuous}),\quad F^{-1}(U)\sim F\\
&\text{Change of variables:} && f_Y(y)=f_X(h(y))\,|h'(y)|,\quad h=g^{-1}\\
&\text{Lognormal:} && f_Y(y)=\tfrac{1}{y\sigma\sqrt{2\pi}}e^{-(\ln y-\mu)^2/(2\sigma^2)},\quad \E[Y^k]=e^{k\mu+k^2\sigma^2/2}\\
&\text{Call payoff:} && \E[(X-K)^+]=\textstyle\int_K^\infty[1-F(x)]\dd x\\
&\quad\text{normal:} && (\mu-K)\Phi(d)+\sigma\phi(d),\quad d=(\mu-K)/\sigma\\
&\quad\text{lognormal:} && e^{\mu+\sigma^2/2}\Phi(d+\sigma)-K\Phi(d),\quad d=(\mu-\ln K)/\sigma\\
&\text{MGFs:} && M_{a+bX}(t)=e^{at}M_X(bt),\quad N(\mu,\sigma^2)\!:e^{\mu t+\sigma^2t^2/2},\quad \chi^2_n\!:(1-2t)^{-n/2}\\
&\text{Portfolio:} && \E[\bm a]\T\bm X=\bm a\T\bm\mu,\quad \Var\bm a\T\bm X=\bm a\T\Sigma\bm a\\
&\text{Conditional normal:} && \E[Y\mid X]=\mu_Y+\rho\tfrac{\tau}{\sigma}(X-\mu_X),\quad \Var(Y\mid X)=(1-\rho^2)\tau^2\\
&\text{PCA:} && \bm p=\Gamma\T(\bm x-\bm\alpha),\ \Cov\bm p=\Lambda,\ \tr\Sigma=\textstyle\sum\lambda_k\\
&\text{Chebyshev, CLT:} && \Prob(|\bar X_n-\mu|>\varepsilon)\le\tfrac{\sigma^2}{n\varepsilon^2},\quad \sqrt n(\bar X_n-\mu)\to N(0,\sigma^2)\\
&\text{Stein:} && \Cov(g(Y),Y^*)=\E[g'(Y)]\Cov(Y,Y^*)\\
&\text{CARA price:} && p_a=\bigl[\E[\tilde F]-A(Q_a+X)\Var\tilde F\bigr]/(1+r_f)\\
&\text{CAPM:} && \E[\tilde r_i]=r_f+\beta_i(\E[\tilde r_m]-r_f),\quad \beta_i=\Cov(\tilde r_i,\tilde r_m)/\Var\tilde r_m
\end{align*}
\end{keyformulas}
